\documentclass[11pt,a4paper]{article}

\usepackage[a4paper,margin=27mm]{geometry}
\usepackage{cmap}
\usepackage[T1]{fontenc}
\usepackage[utf8]{inputenc}
\usepackage{lmodern}
\IfFileExists{glyphtounicode.tex}{\input{glyphtounicode}\pdfgentounicode=1}{}
\usepackage{amsmath,amssymb,amsthm,mathtools}
\numberwithin{equation}{section}
\usepackage{booktabs,longtable,array,tabularx}
\newcolumntype{Y}{>{\raggedright\arraybackslash}X}
\usepackage{graphicx}
\usepackage{enumitem}
\usepackage{microtype}
\usepackage{xcolor}
\usepackage{listings}
\usepackage[hidelinks]{hyperref}

% Anonymous-review defaults.  Replace these macros in the accepted version.
\newcommand{\PaperAuthor}{Anonymous Author(s)}
\newcommand{\PaperAffiliation}{Affiliation withheld for double-anonymous review}
\newcommand{\PaperEmail}{Correspondence details withheld for review}
\newcommand{\PaperORCID}{ORCID withheld for review}
\newcommand{\PaperVersion}{Revision 2.0, 25 July 2026}

\hypersetup{
  pdftitle={Dittert's Conjecture in Dimensions Six through Fifteen: An Exact Computer-Assisted Proof},
  pdfauthor={Anonymous submission},
  pdfsubject={Permanent inequalities and Dittert's conjecture},
  pdfkeywords={permanent, Dittert conjecture, doubly stochastic matrix, Hall condition, stable polynomial, exact Bernstein certificate}
}

\allowdisplaybreaks
\setlength{\parindent}{1.25em}
\setlength{\parskip}{0.15em}

\newtheorem{theorem}{Theorem}[section]
\newtheorem{lemma}[theorem]{Lemma}
\newtheorem{proposition}[theorem]{Proposition}
\newtheorem{corollary}[theorem]{Corollary}
\newtheorem{claim}[theorem]{Claim}
\theoremstyle{definition}
\newtheorem{definition}[theorem]{Definition}
\newtheorem{remark}[theorem]{Remark}

\newcommand{\per}{\operatorname{per}}
\newcommand{\pcap}{\operatorname{Cap}}
\newcommand{\Rge}{\mathbb{R}_{\ge 0}}
\newcommand{\1}{\mathbf{1}}
\newcommand{\cP}{\mathcal{P}}
\newcommand{\cH}{\mathbb{H}}
\newcommand{\e}{\mathrm{e}}
\newcommand{\eps}{\varepsilon}
\newcommand{\supp}{\operatorname{supp}}

\lstdefinestyle{audit}{
  basicstyle=\ttfamily\scriptsize,
  columns=fullflexible,
  keepspaces=true,
  frame=single,
  breaklines=true,
  showstringspaces=false,
  keywordstyle=\bfseries,
  commentstyle=\itshape
}

\title{Dittert's Conjecture in Dimensions Six through Fifteen:\\
An Exact Computer-Assisted Proof}
\author{\PaperAuthor\\[2pt]
\small\PaperAffiliation\\
\small\texttt{\PaperEmail}\\
\small\PaperORCID}
\date{\PaperVersion}

\begin{document}
\maketitle

\begin{abstract}
Let $K_n$ be the simplex of nonnegative $n\times n$ matrices whose entries
sum to $n$.  If $A\in K_n$ has row sums $r_1,\dots,r_n$ and column sums
$c_1,\dots,c_n$, define
\[
 \Phi(A)=\prod_{i=1}^n r_i+\prod_{j=1}^n c_j-\per(A).
\]
Dittert's conjecture asserts that $\Phi$ is uniquely maximized by the
uniform matrix $U_n=n^{-1}J_n$, with maximum $2-n!/n^n$.  We prove the
conjecture for every $6\le n\le15$.  A Hall-deficit dilation converts a
hypothetical boundary maximizer into a doubly stochastic matrix with a forced
zero pattern.  Stable-polynomial capacity yields support-sensitive permanent
bounds and entropy refinements of the sharp rectangular-zero-block
inequality.  A complementary-minor entropy theorem resolves the one-unit
Hall cores in dimension seven.  Dimension six requires a broader
core-marginal entropy estimate, an exact core-scaling stationarity identity,
and a pairwise smoothing reduction to finitely many four-variable rational
polynomials.  Their positivity, together with the remaining scalar
inequalities, is certified by exact rational Bernstein subdivision.  An
independent pairwise-averaging argument proves that every entrywise positive
global maximizer is uniform.  The finite part is verified by explicit rational certificate checkers whose
success conditions remain active under every Python optimization mode.  The
source, independent audit implementations, hashes, and machine-readable
verification manifest are supplied with the paper.
\end{abstract}

\noindent\textbf{Keywords.}
Permanent; Dittert conjecture; doubly stochastic matrix; stable polynomial; polynomial capacity; Hall condition; Birkhoff polytope.

\medskip
\noindent\textbf{2020 Mathematics Subject Classification.}
15A15, 15B51, 05C70.

\section{Introduction}

For an $n\times n$ nonnegative matrix $A=(a_{ij})$, let $r_i$ and $c_j$
denote its row and column sums.  Dittert proposed the extremal problem of
maximizing
\[
 \Phi(A)=\prod_i r_i+\prod_j c_j-\per(A)
\]
on the simplex $K_n$ of matrices whose entries have total sum $n$.  The
conjectured unique maximizer is the uniform matrix $U_n=n^{-1}J_n$.  On the
doubly stochastic polytope the two product terms equal one, so Dittert's
assertion reduces to the van der Waerden permanent theorem of Egorychev and
Falikman \cite{Egorychev1981,Falikman1981}.  The conjecture was recorded by
Minc \cite{Minc1983} and later by Zhan \cite{Zhan2007}.

Sinkhorn proved the case $n=2$ \cite{Sinkhorn1984}, and Hwang proved the case
$n=3$ \cite{Hwang1987}; Hwang had earlier shown that an entrywise positive
global maximizer must be uniform \cite{Hwang1986}.  Knopp and Sinkhorn studied permanent lower bounds for doubly stochastic
matrices with a zero \cite{KnoppSinkhorn1982}, and Cheon and Wanless
developed boundary-exclusion methods and other structural results
\cite{CheonWanless2012}.  Udayan and Somasundaram proved that every
$\Phi$-maximizing matrix on $K_4$ is fully indecomposable
\cite{UdayanSomasundaram2024}.  Recent preprints of Pang and Kafidov claim
the conjecture for $n\ge17$ and $n=16$, respectively
\cite{Pang2026,Kafidov2026}.  These historical statements are not logical
inputs to our proof.  The present paper proves the ten dimensions
\[
 6\le n\le15.
\]

The proof separates the interior and boundary of $K_n$.  The interior
classification is dimension-free: pairwise averaging is exactly flat at a
positive global maximizer, whereas the uniform matrix is a strict local
maximizer.  Cyclic averaging and a backward strictness argument force every
positive global maximizer to be uniform.

For a boundary maximizer, the row- and column-product losses share one
deficit budget.  Relative entropy controls subset sums, and a
max-flow/min-cut criterion shows that a scalar dilation dominates a doubly
stochastic matrix $B$.  At a critical cut, $B$ contains a rectangular zero
block.  Stable-polynomial capacity gives permanent lower bounds adapted to
that support.

Dimensions $8,9,10$ require a second-order scalar comparison and a separate
analysis of one-sided Hall cuts.  Dimensions $11,12,13,14$ are handled by
first-order margins, apart from one deficient-column configuration in
dimension $11$, and dimension $15$ admits a thin--thick dichotomy.
Dimension seven is qualitatively different: three critical cuts with
$n-a-b=1$ survive the scalar method.  A geometric-mean entropy inequality
for complementary permanental minors strengthens the sharp zero-block bound
on these one-unit cores.  Stationarity and conditional product energy then
reduce the remaining cases to three exact two-variable polynomial
alternatives.

Dimension six requires one further level of structure.  Critical cores with
$n-a-b\ge2$ can survive all scalar comparisons.  We prove a
core-marginal entropy refinement of the rectangular-zero-block bound that
records both normalized marginal distributions of the Hall core.  An exact
core-scaling identity couples this entropy to the row and column products.
A pairwise smoothing lemma reduces each side to a uniform distribution on
its support and two levels of row or column sums.  The six Hall block types
then become $98$ explicit four-variable rational polynomials.  One-sided
cuts are excluded by retaining the exact complementary column-product loss.

The finite positivity statements are verified by tensor-product Bernstein
subdivision with exact rational coefficients.  This is a computer-assisted
proof component, not a floating-point experiment: every accepted box has
strictly positive Bernstein coefficients, and every subdivision is exact de
Casteljau averaging.  The dimension-seven-through-fifteen checker uses only
the Python standard library.  The dimension-six generator uses
\texttt{sympy.Rational} only to expand the displayed rational formulas; the
Bernstein transformation, subdivision, and every correctness decision use
exact integers and \texttt{fractions.Fraction}.  Independent implementations
reconstruct the $98$ four-variable certificates, the ten one-variable
certificates, and the three dimension-seven trees.  The complete source code,
a machine-readable manifest, and source hashes reproduce every certificate.

Two ingredients often taken from the literature are proved here rather than
used as black boxes.  We reconstruct the derivative-capacity mechanism
associated with Gurvits \cite{Gurvits2008}, deriving all permanent bounds
needed below, and we give an independent proof of Hwang's positive-maximizer
conclusion.  We make no claim here about the complete status of dimensions
outside $6\le n\le15$: the cited $n=4$, $n=16$, and $n\ge17$ literature has
different publication status and scope, and none of it is needed for our
main theorem.

\subsection*{Proof architecture and certificate scale}
Table~\ref{tab:method-overview} records where analytic structure ends and
finite exact verification begins.
\begin{table}[ht]
\centering
\small
\begin{tabularx}{\textwidth}{@{}c Y Y@{}}
\toprule
Dimensions & Principal analytic mechanism & Exact finite certificate \\
\midrule
$11$--$14$ & first-order Hall margins; one deficient-column refinement for $n=11$ & $35,42,49,56$ rational cases \\
$15$ & thin--thick support dichotomy & two rational margins and six integer inequalities \\
$8$--$10$ & second-order Hall margins and one-sided entropy & $12,16,20$ rational Hall cases plus scalar certificates \\
$7$ & complementary-minor entropy and single-bridge stationarity & three bivariate Bernstein trees \\
$6$ & full core-marginal entropy, core stationarity, and two-level smoothing & $98$ four-variable and ten univariate Bernstein certificates \\
\bottomrule
\end{tabularx}
\caption{Dimension-by-dimension analytic mechanism and exact certificate size.}
\label{tab:method-overview}
\end{table}

The logical dependency is acyclic.  Its boundary branch is
\[
 \boxed{\text{stable capacity}}
 \Longrightarrow
 \boxed{\text{support/entropy permanent bounds}}
 \Longrightarrow
 \boxed{\text{Hall-core bounds}},
\]
\[
 \boxed{\text{Hall-core bounds}}
 \Longrightarrow
 \boxed{\text{boundary exclusion}}.
\]
The interior branch is
\[
 \boxed{\text{pairwise flatness and averaging convergence}}
 \Longrightarrow
 \boxed{\text{positive maximizers converge to }U_n},
\]
with strict local maximality providing the backward-uniqueness step.
The main theorem combines these two branches only in
Section~\ref{sec:completion}.

\section{Notation and the main theorem}

Let $S_n$ be the symmetric group.  For $A=(a_{ij})\in\Rge^{n\times n}$,
\[
 \per(A)=\sum_{\sigma\in S_n}\prod_{i=1}^n a_{i,\sigma(i)}.
\]
Let
\[
 K_n=\left\{A\in\Rge^{n\times n}:\sum_{i,j=1}^n a_{ij}=n\right\}.
\]
For $A\in K_n$, put
\[
 r_i=\sum_{j=1}^n a_{ij},\qquad c_j=\sum_{i=1}^n a_{ij},
\]
and
\[
 \Phi(A)=\prod_{i=1}^n r_i+\prod_{j=1}^n c_j-\per(A).
\]
Let $J_n=\1\1^{\mathsf T}$, let
\[
 U_n=\frac1nJ_n,
\]
and define
\[
 \gamma_n=\per(U_n)=\frac{n!}{n^n}.
\]
A matrix is called \emph{positive} if every entry is strictly positive.  A matrix is \emph{doubly stochastic} if every row sum and every column sum equals one.
Symbols introduced inside a dimension-specific subsection are local to that
subsection; repeated letters such as $H$, $K$, and $q$ are always redefined
before reuse.  The global notation above remains fixed throughout.

\begin{theorem}[Main theorem]\label{thm:main}
For every integer $n$ with $6\le n\le15$ and every $A\in K_n$,
\[
 \Phi(A)\le 2-\frac{n!}{n^n}.
\]
Equality holds if and only if $A=U_n$.
\end{theorem}

The simplex $K_n$ is compact and $\Phi$ is continuous, so global maximizers exist.  We first develop two pieces of infrastructure that are valid in arbitrary dimension: permanent bounds from capacity and the classification of positive maximizers.

\section{Stable-polynomial capacity and permanent bounds}\label{sec:capacity}

\subsection{Stability and capacity}

Let
\[
 \cH=\{z\in\mathbb C:\operatorname{Im}z>0\}.
\]
A polynomial $p\in\mathbb R[x_1,\dots,x_m]$ is \emph{real stable} if
\[
 p(z_1,\dots,z_m)\ne0
 \qquad(z_1,\dots,z_m\in\cH).
\]
For a homogeneous polynomial $p$ of degree $m$ in $m$ variables with
nonnegative coefficients, define
\[
 \pcap(p)=\inf_{x_1,\dots,x_m>0}
 \frac{p(x_1,\dots,x_m)}{x_1\cdots x_m}.
\]
We also set $\pcap(0)=0$.  If a nonzero polynomial in this quotient is
independent of some denominator variable $x_j$, then its capacity is zero:
fix the remaining variables and let $x_j\to\infty$.  We shall use this
observation to verify explicitly that every variable extracted below has
positive degree.
Set
\[
 G(1)=1,
 \qquad
 G(k)=\left(\frac{k-1}{k}\right)^{k-1}\quad(k\ge2),
\]
and
\[
 v_m=\frac{m!}{m^m}\quad(m\ge1),
 \qquad v_0=1.
\]
The telescoping identity
\begin{equation}\label{eq:telescoping}
 \prod_{k=1}^mG(k)=v_m
\end{equation}
will be used repeatedly.  The sequence $G(k)$ is decreasing.  Indeed,
$G(k)^{-1}=(1+1/(k-1))^{k-1}$, and the function
\[
 x\longmapsto x\log(1+1/x)
\]
is increasing on $(0,\infty)$ because
$\log(1+y)>y/(1+y)$ for $y>0$.

We use two classical closure properties of real stability.  They are included for completeness.

\begin{lemma}[Closure under specialization and differentiation]\label{lem:stability-closure}
Let $p$ be real stable.
\begin{enumerate}[label=\textup{(\roman*)}]
\item Specializing one variable to a real number produces either a real-stable polynomial or the zero polynomial.
\item Every nonzero partial derivative of $p$ is real stable.
\end{enumerate}
\end{lemma}

\begin{proof}
For real specialization, fix $a\in\mathbb R$ and replace it first by $a+i\eps$.  The resulting polynomials are zero-free on the product of upper half-planes and converge locally uniformly as $\eps\downarrow0$.  Hurwitz's theorem implies that the limit is either zero-free there or identically zero.

For differentiation, write $p(z',w)=\sum_{r=0}^d a_r(z')w^r$ with $a_d$ not identically zero.  The functions $(iR)^{-d}p(z',iR)$ are zero-free for $z'\in\cH^{m-1}$ and converge locally uniformly to $a_d(z')$ as $R\to\infty$.  Hurwitz's theorem shows that $a_d$ is real stable, so $a_d(z')\ne0$ for every $z'\in\cH^{m-1}$.  Thus, after fixing such a $z'$, the univariate polynomial $w\mapsto p(z',w)$ has exact degree $d$, and every one of its roots lies in the closed lower half-plane.  By the Gauss--Lucas theorem, the roots of its derivative lie in the convex hull of those roots and hence also in the closed lower half-plane.  Therefore the derivative does not vanish for $w\in\cH$.  Since $z'$ was arbitrary, the partial derivative is real stable.  The zero-polynomial case is excluded by hypothesis.
\end{proof}

The next one-variable estimate is the quantitative core of the capacity method.

\begin{lemma}[One-variable capacity estimate]\label{lem:one-variable}
Let $f$ be a nonzero univariate polynomial with nonnegative coefficients.  Suppose all zeros of $f$ are real and nonpositive, and let $d=\deg f\ge1$.  Then
\begin{equation}\label{eq:one-variable}
 f'(0)\ge G(d)\inf_{t>0}\frac{f(t)}t.
\end{equation}
\end{lemma}

\begin{proof}
Assume first that $f(0)>0$.  Then
\[
 f(t)=f(0)\prod_{\ell=1}^d(1+\lambda_\ell t),
 \qquad \lambda_\ell\ge0.
\]
Let
\[
 s=\sum_{\ell=1}^d\lambda_\ell=\frac{f'(0)}{f(0)}.
\]
Since $d\ge1$ is the exact degree, $s>0$.  By AM--GM,
\[
 \prod_{\ell=1}^d(1+\lambda_\ell t)
 \le \left(1+\frac{st}{d}\right)^d.
\]
If $d\ge2$, choose $t=d/[s(d-1)]$.  This gives
\[
 \inf_{t>0}\frac{f(t)}t
 \le \frac{f'(0)}{G(d)}.
\]
For $d=1$, the same conclusion follows by letting $t\to\infty$.

If $f(0)=0$, write $f(t)=t^r g(t)$ with $g(0)>0$.  For $r=1$,
\[
 f'(0)=g(0)\ge G(d)\inf_{t>0}\frac{f(t)}t,
\]
because the infimum is at most the limit $g(0)$ as $t\downarrow0$.  If $r\ge2$, then both $f'(0)$ and the infimum in \eqref{eq:one-variable} are zero.
\end{proof}

\begin{theorem}[Derivative-capacity inequality]\label{thm:derivative-capacity}
Let $p$ be a homogeneous real-stable polynomial of degree $m$ in $m$ variables, with nonnegative coefficients.  Put
\[
 d=\deg_{x_m}p
\]
and
\[
 q(x_1,\dots,x_{m-1})
 =\left.\frac{\partial p}{\partial x_m}\right|_{x_m=0}.
\]
If $d\ge1$, then
\begin{equation}\label{eq:derivative-capacity}
 \pcap(q)\ge G(d)\pcap(p).
\end{equation}
\end{theorem}

\begin{proof}
If $q=0$, then for every $x'>0$ the polynomial $f(t)=p(x',t)$ satisfies
$f'(0)=0$.  Lemma~\ref{lem:one-variable} therefore gives
$\inf_{t>0}f(t)/t=0$ for every $x'$, and hence $\pcap(p)=0$; the assertion
follows from the convention $\pcap(0)=0$.  We may consequently assume
$q\ne0$.

Fix $x'=(x_1,\dots,x_{m-1})>0$ and set
\[
 f(t)=p(x',t).
\]
By Lemma~\ref{lem:stability-closure}, $f$ is stable as a univariate real polynomial.  Hence all its zeros are real: a nonreal zero below the real axis would have a conjugate zero above it.  Nonnegative coefficients rule out positive zeros.  Thus Lemma~\ref{lem:one-variable} applies.

For every $t>0$, the definition of capacity gives
\[
 \frac{f(t)}{t}\ge \pcap(p)x_1\cdots x_{m-1}.
\]
Therefore
\[
 q(x')=f'(0)
 \ge G(d)\pcap(p)x_1\cdots x_{m-1}.
\]
Divide by $x_1\cdots x_{m-1}$ and take the infimum over $x'>0$.
\end{proof}

\begin{corollary}[Square-free coefficient extraction]\label{cor:squarefree}
Let $p$ be homogeneous, real stable, of degree $m$ in $m$ variables, and have nonnegative coefficients.  Then the coefficient of $x_1\cdots x_m$ is at least
\[
 v_m\pcap(p).
\]
\end{corollary}

\begin{proof}
If $\pcap(p)=0$, the assertion is immediate.  Otherwise differentiate successively in $x_m,x_{m-1},\dots,x_1$, setting each differentiated variable to zero.  Theorem~\ref{thm:derivative-capacity} keeps every intermediate capacity positive.  At the stage with $j$ variables remaining, the relevant degree is at most $j$ and at least one; if it were zero, the polynomial would be independent of a denominator variable and would have capacity zero.  Since $G$ is decreasing, Theorem~\ref{thm:derivative-capacity} contributes at least $G(j)$.  Equation~\eqref{eq:telescoping} completes the proof.
\end{proof}

\subsection{The support-sensitive permanent inequality}

\begin{theorem}[Support-sensitive permanent bound]\label{thm:support-bound}
Let $B=(b_{ij})$ be an $n\times n$ doubly stochastic matrix.  Order its columns arbitrarily, and let $d_j$ be the number of positive entries in the $j$th ordered column.  Then
\begin{equation}\label{eq:support-bound}
 \per(B)\ge
 \prod_{j=1}^nG\bigl(\min\{j,d_j\}\bigr).
\end{equation}
\end{theorem}

\begin{proof}
Consider
\[
 P_B(x_1,\dots,x_n)
 =\prod_{i=1}^n\left(\sum_{j=1}^n b_{ij}x_j\right).
\]
Every linear factor has positive imaginary part when all variables lie in $\cH$, so $P_B$ is real stable.  Weighted AM--GM gives, for $x_j>0$,
\[
 \sum_jb_{ij}x_j\ge\prod_jx_j^{b_{ij}}.
\]
Multiplying over the rows and using the column sums,
\[
 P_B(x)\ge x_1\cdots x_n.
\]
Equality holds at $x_1=\cdots=x_n=1$, and therefore
\begin{equation}\label{eq:PB-capacity}
 \pcap(P_B)=1.
\end{equation}

Set $p_n=P_B$ and recursively
\[
 p_{j-1}
 =\left.\frac{\partial p_j}{\partial x_j}\right|_{x_j=0}
 \qquad(j=n,n-1,\dots,1).
\]
At stage $j$, the polynomial has degree $j$ in $j$ variables.
Differentiation in other variables cannot increase its degree in $x_j$, so
\[
 \deg_{x_j}p_j\le\min\{j,d_j\}.
\]
Every intermediate capacity is positive.  Therefore $p_j$ cannot be
independent of $x_j$, for otherwise the denominator-variable observation
above would give $\pcap(p_j)=0$.  Thus
$1\le\deg_{x_j}p_j\le\min\{j,d_j\}$, and all hypotheses of
Theorem~\ref{thm:derivative-capacity} are met.  Using the monotonicity of
$G$ and \eqref{eq:PB-capacity},
\[
 p_0\ge\prod_{j=1}^nG\bigl(\min\{j,d_j\}\bigr).
\]
The scalar $p_0$ is the coefficient of $x_1\cdots x_n$ in $P_B$, namely $\per(B)$.
\end{proof}

\subsection{Zero-pattern and entropy consequences}

\begin{corollary}[One zero]\label{cor:one-zero}
If a doubly stochastic $B$ has at least one zero entry, then
\begin{equation}\label{eq:kappa}
 \per(B)\ge\kappa_n,
 \qquad
 \kappa_n:=v_{n-1}G(n-1)
 =\frac{(n-2)!(n-2)^{n-2}}{(n-1)^{2n-4}}.
\end{equation}
Moreover $\kappa_n>\gamma_n$ for every $n\ge3$.
\end{corollary}

\begin{proof}
Put a column containing a zero last in Theorem~\ref{thm:support-bound}.  Its support size is at most $n-1$, giving \eqref{eq:kappa}.

For $m=n-1\ge2$,
\[
 \frac{\kappa_n}{\gamma_n}
 =\left(1+\frac1m\right)^m
  \left(1-\frac1m\right)^{m-1}
 =\left(1+\frac1m\right)
  \left(1-\frac1{m^2}\right)^{m-1}.
\]
Bernoulli's inequality yields
\[
 \left(1-\frac1{m^2}\right)^{m-1}
 \ge1-\frac{m-1}{m^2}
 =1-\frac1m+\frac1{m^2}.
\]
Multiplication by $1+1/m$ gives $1+1/m^3>1$.
\end{proof}

\begin{theorem}[Best-possible rectangular-zero-block bound]\label{thm:rectangular}
Let $B$ be doubly stochastic and suppose that, after row and column permutations, it contains an $a\times b$ zero block, where
\[
 a,b\ge1,
 \qquad
 p:=n-a-b\ge1.
\]
Then
\begin{equation}\label{eq:sharp-block}
 \per(B)\ge
 L^{\sharp}_{n,a,b}:=
 \frac{v_{n-a}v_{n-b}}{v_p}.
\end{equation}
The constant is best possible; equality is attained by the block-constant
matrix displayed in the proof.  No classification of all equality cases is
asserted here.
\end{theorem}

\begin{proof}
Let $C$ be the $b$ columns in the zero block and let $R$ be its $a$ rows.  In $P_B$, every variable indexed by $C$ occurs only in the $n-a$ row factors indexed by $R^c$.

Differentiate successively in the $b$ variables indexed by $C$, setting
each differentiated variable to zero.  After $k$ such differentiations,
every surviving product-rule term has differentiated $k$ distinct row
factors: each row factor is linear and cannot be differentiated twice.
Consequently, the degree of the next special variable is at most $n-a-k$.
The current capacity remains positive, so that degree is also at least one;
otherwise the current polynomial would omit a denominator variable and have
zero capacity.  Theorem~\ref{thm:derivative-capacity} therefore supplies the
factor
\[
 \prod_{k=0}^{b-1}G(n-a-k)
 =\prod_{j=p+1}^{n-a}G(j)
 =\frac{v_{n-a}}{v_p}.
\]
A homogeneous stable polynomial of degree $n-b$ in $n-b$ variables remains.  Corollary~\ref{cor:squarefree} supplies the additional factor $v_{n-b}$.  Since $\pcap(P_B)=1$, \eqref{eq:sharp-block} follows.

To prove best-possible attainment, consider
\[
 B_0=
 \begin{pmatrix}
  0_{a\times b} & \dfrac1{n-b}J_{a,n-b}\\[2mm]
  \dfrac1{n-a}J_{n-a,b} &
  \dfrac{p}{(n-a)(n-b)}J_{n-a,n-b}
 \end{pmatrix}.
\]
It is doubly stochastic.  Every contributing permutation uses $a$ upper-right entries, $b$ lower-left entries, and $p$ lower-right entries.  Their number is
\[
 \frac{(n-b)!}{p!}\cdot\frac{(n-a)!}{p!}\cdot p!
 =\frac{(n-a)!(n-b)!}{p!}.
\]
Multiplying by the common weight gives
\[
 \per(B_0)
 =\frac{(n-a)!(n-b)!}{p!}
  \frac{p^p}{(n-a)^{n-a}(n-b)^{n-b}}
 =\frac{v_{n-a}v_{n-b}}{v_p}.
\]
\end{proof}

For the middle dimensions it is convenient to record a weaker consequence with a simpler form.

\begin{corollary}\label{cor:weak-block}
Suppose $1\le a\le b$ and a doubly stochastic $B$ contains an $a\times b$ zero block.  Then
\begin{equation}\label{eq:weak-block}
 \per(B)\ge L_{n,a,b}:=v_{n-a}G(n-b)^a.
\end{equation}
In particular, a $2\times2$ zero block gives
\begin{equation}\label{eq:two-by-two}
 \per(B)\ge v_{n-2}G(n-2)^2.
\end{equation}
\end{corollary}

\begin{proof}
Transpose first.  The transposed zero block has $a$ columns, each supported on at most $n-b$ rows.  Put those $a$ columns last in Theorem~\ref{thm:support-bound}; the first $n-a$ factors give $v_{n-a}$ and the last $a$ factors give $G(n-b)^a$.
\end{proof}

\begin{theorem}[Column-entropy permanent bound]\label{thm:column-entropy}
Let $B$ be doubly stochastic, and let $b=(b_1,\dots,b_n)$ be any one of its columns.  Put
\[
 s(b)=\sum_{i=1}^n\left(b_i-\frac1n\right)^2.
\]
Then
\begin{equation}\label{eq:entropy-permanent}
 \per(B)\ge\gamma_n\exp\left(\frac{s(b)}2\right).
\end{equation}
\end{theorem}

\begin{proof}
It is enough first to consider $0<b_i<1$.  For a general doubly stochastic
matrix, apply the positive-case argument to
$B_\eps=(1-\eps)B+\eps U_n$.  The entropy expression is continuous under
$0\log0=0$, and the permanent is polynomial in the entries, so the asserted
inequality follows by letting $\eps\downarrow0$.  Take the selected column
to be the last one and set
\[
 L_i(x_1,\dots,x_{n-1})=\sum_{j<n}b_{ij}x_j.
\]
The derivative polynomial is
\[
 q=\left.\frac{\partial P_B}{\partial x_n}\right|_{x_n=0}
 =\sum_i b_i\prod_{k\ne i}L_k.
\]
Weighted AM--GM, with weights $b_i$, gives
\[
 q\ge\prod_iL_i^{1-b_i}.
\]
A second weighted AM--GM inequality gives
\[
 L_i\ge(1-b_i)
 \prod_{j<n}x_j^{b_{ij}/(1-b_i)}.
\]
Therefore
\[
 q(x)\ge
 \left(\prod_i(1-b_i)^{1-b_i}\right)x_1\cdots x_{n-1},
\]
so
\[
 \pcap(q)\ge\prod_i(1-b_i)^{1-b_i}.
\]
Corollary~\ref{cor:squarefree} yields
\begin{equation}\label{eq:entropy-intermediate}
 \per(B)\ge
 v_{n-1}\prod_i(1-b_i)^{1-b_i}.
\end{equation}

Let $h(t)=(1-t)\log(1-t)$.  Since $h''(t)=1/(1-t)\ge1$ on $[0,1)$,
\[
 h(t)\ge h(1/n)+h'(1/n)(t-1/n)+\frac12(t-1/n)^2.
\]
Summing at $t=b_i$, the linear terms cancel.  Hence
\[
 \sum_i h(b_i)
 \ge nh(1/n)+\frac12s(b).
\]
Now $\exp(nh(1/n))=G(n)$ and $v_{n-1}G(n)=\gamma_n$.  Substitution into \eqref{eq:entropy-intermediate} proves \eqref{eq:entropy-permanent}.
\end{proof}


\subsection{Complement entropy and core-marginal refinements}

The column-entropy estimate retains one column of a doubly stochastic
matrix.  The lower-dimensional boundary arguments also require the dual
information carried by all complementary minors of a column and, in
Dimension~6, by the marginal distributions of a whole Hall core.

For a probability vector $q=(q_1,\dots,q_m)$, put
\begin{equation}\label{eq:Psi-definition}
 \Psi_m(q)=
 \sum_{i=1}^m(1-q_i)\log(1-q_i)
 -m\left(1-\frac1m\right)\log\left(1-\frac1m\right),
\end{equation}
with the convention $0\log0=0$, and write
\[
 s_m(q)=\sum_{i=1}^m\left(q_i-\frac1m\right)^2.
\]

\begin{lemma}[Quadratic complement-entropy bounds]\label{lem:Psi-bounds}
For every probability vector $q$ of length $m\ge2$,
\begin{equation}\label{eq:Psi-bounds}
 \Psi_m(q)\ge
 \left(\frac12+\frac1{3m}\right)s_m(q)
 \ge\frac12s_m(q).
\end{equation}
\end{lemma}

\begin{proof}
Let $h(x)=(1-x)\log(1-x)$.  Since
\[
 h''(x)=\frac1{1-x}\ge1+x
 \qquad(0\le x<1),
\]
For $q_i<1$, Taylor's formula with integral remainder about $1/m$ gives, for
$d_i=q_i-1/m$,
\[
 h(q_i)\ge h(1/m)+h'(1/m)d_i
 +\frac{1+1/m}{2}d_i^2+\frac16d_i^3.
\]
Because $d_i\ge-1/m$, one has $d_i^3\ge-m^{-1}d_i^2$.
Summation and $\sum_i d_i=0$ prove the first inequality.  Boundary
coordinates $q_i=1$ follow by continuity.  The second inequality is
immediate.
\end{proof}

\begin{theorem}[Complementary-minor entropy]\label{thm:minor-entropy}
Let $M$ be an $m\times m$ doubly stochastic matrix, let
$q=(q_1,\dots,q_m)$ be one of its columns, and let
\[
 d_i=\per M(i\mid j_0)
\]
be the permanent of the submatrix obtained by deleting row $i$ and the
selected column $j_0$.  With the convention that a factor of exponent zero
is omitted,
\begin{equation}\label{eq:minor-entropy-exact}
 \prod_{i=1}^m d_i^{q_i}
 \ge v_m\exp\bigl(\Psi_m(q)\bigr).
\end{equation}
Consequently,
\begin{equation}\label{eq:minor-entropy}
 \prod_{i=1}^m d_i^{q_i}
 \ge v_m\exp\left\{
 \left(\frac12+\frac1{3m}\right)s_m(q)
 \right\}
 \ge v_m\exp\left(\frac{s_m(q)}2\right).
\end{equation}
\end{theorem}

\begin{proof}
It is enough first to treat the positive case.  For a general $M$, apply
the argument to $M_\eps=(1-\eps)M+\eps U_m$ and let $\eps\downarrow0$.
All selected-column coordinates and complementary minors are positive for
$\eps>0$.  If $q_i(0)=0$, then $q_i(\eps)=O(\eps)$, while
$d_i(\eps)$ is a nonzero polynomial in $\eps$; hence
$q_i(\eps)\log d_i(\eps)=O(\eps|\log\eps|)$.  If $q_i(0)>0$, then
$d_i(0)>0$: otherwise the left side of the positive-$\eps$ inequality would
tend to zero, whereas its right side tends to the strictly positive number
$v_m\exp(\Psi_m(q(0)))$.  Thus every logarithmic term has the asserted
limit, and the product inequality passes to $\eps=0$.

Take the selected column to be the last one and write
\[
 L_i(x_1,\dots,x_{m-1})=\sum_{j<m}m_{ij}x_j.
\]
For arbitrary $t_1,\dots,t_m>0$, set
\[
 q_t(x)=\sum_{i=1}^m q_i t_i\prod_{k\ne i}L_k(x).
\]
This is homogeneous, real stable, and of degree $m-1$, because
\[
 q_t=
 \left.\frac{\partial}{\partial z}
 \prod_{i=1}^m\bigl(L_i(x)+q_it_i z\bigr)\right|_{z=0}.
\]
Weighted AM--GM, first with weights $q_i$ and then within each row, gives
\[
 \begin{aligned}
 q_t
 &=\left(\prod_iL_i\right)\sum_iq_i\frac{t_i}{L_i}\\
 &\ge\left(\prod_i t_i^{q_i}\right)
      \prod_iL_i^{1-q_i}\\
 &\ge\left(\prod_i t_i^{q_i}\right)
      \left(\prod_i(1-q_i)^{1-q_i}\right)
      x_1\cdots x_{m-1}.
 \end{aligned}
\]
Thus
\[
 \pcap(q_t)\ge
 \prod_i t_i^{q_i}\prod_i(1-q_i)^{1-q_i}.
\]
The coefficient of $x_1\cdots x_{m-1}$ in $q_t$ is
$\sum_iq_it_id_i$, so Corollary~\ref{cor:squarefree} yields
\[
 \sum_iq_it_id_i\ge
 v_{m-1}\prod_i t_i^{q_i}\prod_i(1-q_i)^{1-q_i}.
\]
This inequality implies $d_i>0$ whenever $q_i>0$.  Put $t_i=d_i^{-1}$
on the support of $q$.  Since $\sum_iq_i=1$,
\[
 \prod_i d_i^{q_i}
 \ge v_{m-1}\prod_i(1-q_i)^{1-q_i}
 =v_m\exp\bigl(\Psi_m(q)\bigr),
\]
where $v_{m-1}G(m)=v_m$.  Lemma~\ref{lem:Psi-bounds} gives the stated
consequences.
\end{proof}

The case in which a critical zero block leaves only one unit of mass in the
opposite corner will be called a \emph{single-bridge} configuration.

\begin{theorem}[Single-bridge zero-block entropy]\label{thm:single-bridge-entropy}
Let $B$ be an $n\times n$ doubly stochastic matrix with a zero block
$B[R,C]=0$, where $|R|=a$, $|C|=b$, and $a+b=n-1$.
Put $I=R^c$, $J=C^c$, $m=b+1$, and $\ell=a+1$.
The core $X=B[I,J]$ has total mass one.  Let $\xi$ and $\eta$ be its row-
and column-sum vectors, and put
\[
 s_I=s_m(\xi),\qquad s_J=s_\ell(\eta).
\]
Then
\begin{equation}\label{eq:single-bridge-entropy-exact}
 \per(B)\ge
 v_mv_\ell\exp\bigl(\Psi_m(\xi)+\Psi_\ell(\eta)\bigr).
\end{equation}
In particular, with $c_q=\frac12+\frac1{3q}$,
\begin{equation}\label{eq:single-bridge-entropy-strong}
 \per(B)\ge v_mv_\ell\exp(c_ms_I+c_\ell s_J)
 \ge v_mv_\ell\exp\left(\frac{s_I+s_J}{2}\right).
\end{equation}
The constant $v_mv_\ell$ is $L^{\sharp}_{n,a,b}$.
\end{theorem}

\begin{proof}
Every perfect matching uses exactly one edge of $X$.  For $i\in I$ and
$j\in J$, put
\[
 e_i=\per B[I\setminus\{i\},C],
 \qquad
 d_j=\per B[R,J\setminus\{j\}].
\]
Expansion according to the unique core edge gives
\[
 \per(B)=\sum_{i\in I,\,j\in J}x_{ij}e_id_j.
\]
The $x_{ij}$ sum to one, so weighted AM--GM implies
\[
 \per(B)\ge
 \prod_{i\in I}e_i^{\xi_i}
 \prod_{j\in J}d_j^{\eta_j}.
\]
Adjoin $\xi$ as a last column to $B[I,C]$.  The resulting $m\times m$
matrix is doubly stochastic and its complementary minors are the $e_i$.
Similarly, adjoining $\eta$ to the transpose of $B[R,J]$ gives an
$\ell\times\ell$ doubly stochastic matrix whose complementary minors are
the $d_j$.  Apply Theorem~\ref{thm:minor-entropy} on both sides, and then
Lemma~\ref{lem:Psi-bounds}.
\end{proof}

The next estimate retains one supported column inside an arbitrary
rectangular zero block.

\begin{theorem}[Distinguished zero-block column]\label{thm:distinguished-column}
Let $B$ be an $n\times n$ doubly stochastic matrix containing an
$a\times b$ zero block $B[R,C]=0$, where
\[
 a,b\ge1,\qquad p=n-a-b\ge1.
\]
Put $I=R^c$ and $m=n-a$.  Select a column $j_0\in C$, and let
$q=(b_{ij_0})_{i\in I}$ be its probability vector on $I$.  Then
\begin{equation}\label{eq:distinguished-column}
 \per(B)\ge
 L^{\sharp}_{n,a,b}\exp\bigl(\Psi_m(q)\bigr).
\end{equation}
Consequently,
\[
 \per(B)\ge
 L^{\sharp}_{n,a,b}\exp\left(\frac{s_m(q)}2\right).
\]
The transposed assertion holds for a selected row in $R$.
\end{theorem}

\begin{proof}
First assume that $B$ is positive on every position outside the prescribed
zero block.  Differentiate the product polynomial $P_B$ first in the
selected column variable and set that variable to zero.  With
\[
 L_i(x)=\sum_{j\ne j_0}b_{ij}x_j,
\]
the resulting homogeneous real-stable polynomial is
\[
 Q(x)=\sum_{i\in I}q_i\prod_{k\ne i}L_k(x).
\]
Exactly as in the column-entropy proof, two weighted AM--GM inequalities give
\[
 \pcap(Q)\ge\prod_{i\in I}(1-q_i)^{1-q_i}.
\]
Differentiate next in the remaining $b-1$ variables indexed by $C$.
Their successive degrees are at most
\[
 m-1,m-2,\dots,p+1.
\]
At every stage the current capacity is positive, so the variable being
extracted has degree at least one.  Thus the derivative-capacity theorem is
applicable at each step and supplies
$v_{m-1}/v_p$.  Square-free coefficient extraction in the remaining
$n-b$ variables supplies $v_{n-b}$.  Therefore
\[
 \per(B)\ge
 \frac{v_{m-1}v_{n-b}}{v_p}
 \prod_{i\in I}(1-q_i)^{1-q_i}.
\]
Since $v_{m-1}G(m)=v_m$, this is exactly
\eqref{eq:distinguished-column}.

For a general $B$, mix it with the block-constant attaining matrix from
Theorem~\ref{thm:rectangular}.  This preserves the prescribed zero block and
is positive on every allowed entry.  Let the mixing parameter tend to zero;
the permanent is polynomial in the entries and the entropy term is continuous
under the convention $0\log0=0$.
\end{proof}

We shall use the distinguished-column estimate through the following
symmetric consequence.  Let $B[R,C]=0$ be as above, put
\[
 I=R^c,\quad J=C^c,\quad m=n-a,\quad \ell=n-b,
\]
and note that the core $B[I,J]$ has total mass $p$.  Define its normalized
row and column marginals by
\[
 \xi_i=\frac1p\sum_{j\in J}b_{ij},
 \qquad
 \eta_j=\frac1p\sum_{i\in I}b_{ij},
\]
and set $S_\xi=s_m(\xi)$ and $S_\eta=s_\ell(\eta)$.

\begin{proposition}[Core-marginal zero-block entropy]
\label{prop:core-marginal-entropy}
Let $L=L^{\sharp}_{n,a,b}$.  If $p\ge2$, then
\begin{equation}\label{eq:core-entropy-general}
 \per(B)\ge
 L\left(1+c_{a,b}(S_\xi+S_\eta)\right),
 \qquad
 c_{a,b}=\frac{p^2}{2(a^2+b^2)}.
\end{equation}
If $p=1$, then
\begin{equation}\label{eq:core-entropy-p1}
 \per(B)\ge
 L\left(1+c_mS_\xi+c_\ell S_\eta\right),
 \qquad c_q=\frac12+\frac1{3q}.
\end{equation}
\end{proposition}

\begin{proof}
Assume first that $p\ge2$.  For each of the $b$ columns indexed by $C$,
view its entries on $I$ as a probability vector.  Their average is
\[
 \bar q_i=\frac{1-p\xi_i}{b},
\]
and hence
\[
 \sum_{i\in I}\left(\bar q_i-\frac1m\right)^2
 =\frac{p^2}{b^2}S_\xi.
\]
Theorem~\ref{thm:distinguished-column}, convexity of $\Psi_m$, and
Lemma~\ref{lem:Psi-bounds} give
\[
 \log\frac{\per(B)}L
 \ge\max_{j\in C}\Psi_m(q^{(j)})
 \ge\frac1b\sum_{j\in C}\Psi_m(q^{(j)})
 \ge\Psi_m(\bar q)
 \ge\frac{p^2}{2b^2}S_\xi.
\]
Transposition gives
\[
 \log\frac{\per(B)}L\ge\frac{p^2}{2a^2}S_\eta.
\]
The maximum of these two lower bounds is at least their weighted mean with
weights $b^2/(a^2+b^2)$ and $a^2/(a^2+b^2)$.  Thus
\[
 \log\frac{\per(B)}L
 \ge\frac{p^2}{2(a^2+b^2)}(S_\xi+S_\eta).
\]
Use $\exp x\ge1+x$.

If $p=1$, apply Theorem~\ref{thm:single-bridge-entropy} and then
$\exp x\ge1+x$.
\end{proof}



\section{Positive global maximizers are uniform}\label{sec:positive}

This section gives an independent proof of the interior classification.  No permanent lower bound is used.

\subsection{Flatness under pairwise averaging}

Choose two columns $u,v$ of a positive matrix $A\in K_n$.  For $t\in\mathbb R$, replace them by
\begin{equation}\label{eq:pair-path}
 u(t)=(1-t)u+tv,
 \qquad
 v(t)=tu+(1-t)v,
\end{equation}
and leave every other column fixed.  Denote the resulting matrix by $A(t)$.

\begin{lemma}[Pairwise averaging is flat at a positive maximizer]\label{lem:flat-average}
If $A$ is a positive global maximizer of $\Phi$, then
\begin{equation}\label{eq:flat-path}
 \Phi(A(t))=\Phi(A)
 \qquad\text{for every }t\in\mathbb R.
\end{equation}
In particular, replacing two columns by their arithmetic mean preserves the global maximum value.  The same statement holds for two rows.
\end{lemma}

\begin{proof}
Because $A$ is positive, $A(t)\in K_n$ for all $t$ in an open interval containing zero.  The row sums are unchanged along the path.  The product of the column sums is a quadratic polynomial in $t$, and multilinearity of the permanent in the two changed columns shows that $\per(A(t))$ is also quadratic.  Thus
\[
 F(t):=\Phi(A(t))
\]
has degree at most two.

Replacing $t$ by $1-t$ interchanges the two columns, so $F(t)=F(1-t)$.  Hence
\[
 F(t)=F(0)+\lambda t(1-t)
\]
for some real $\lambda$.  Since $A=A(0)$ is a local maximum along a two-sided feasible interval, $F'(0)=0$, and therefore $\lambda=0$.  The row version follows by transposition.
\end{proof}

\subsection{Strict local maximality of the uniform matrix}

\begin{lemma}[Hessian at the uniform matrix]\label{lem:hessian}
Let $X=(x_{ij})$ satisfy
\[
 \sum_{i,j}x_{ij}=0.
\]
Put
\[
 u_i=\sum_jx_{ij},
 \qquad
 v_j=\sum_ix_{ij},
\]
and
\[
 Q=\sum_{i,j}x_{ij}^2,
 \qquad
 R_2=\sum_i u_i^2,
 \qquad
 C_2=\sum_jv_j^2.
\]
Then
\begin{equation}\label{eq:hessian}
 \left.\frac{d^2}{dt^2}\Phi(U_n+tX)\right|_{t=0}
 =-(1-c_n)(R_2+C_2)-c_nQ,
\end{equation}
where
\[
 c_n=\frac{(n-2)!}{n^{n-2}}.
\]
Consequently, the Hessian of $\Phi$ at $U_n$ is negative definite on the tangent hyperplane of $K_n$.
\end{lemma}

\begin{proof}
The row sums of $U_n+tX$ are $1+tu_i$.  Since $\sum_i u_i=0$,
\[
 \left.\frac{d^2}{dt^2}\prod_i(1+tu_i)\right|_{t=0}
 =-R_2.
\]
The column-product contribution is similarly $-C_2$.

For the permanent, direct differentiation gives
\begin{equation}\label{eq:per-second}
 \left.\frac{d^2}{dt^2}\per(U_n+tX)\right|_{t=0}
 =c_n\sum_{\substack{i\ne k\\j\ne\ell}}x_{ij}x_{k\ell}.
\end{equation}
Indeed, for fixed ordered pairs $i\ne k$ and $j\ne\ell$, there are $(n-2)!$ permutations satisfying $\sigma(i)=j$ and $\sigma(k)=\ell$, while the remaining $n-2$ entries contribute $n^{-(n-2)}$.

Expanding the identity
\[
 0=\left(\sum_{i,j}x_{ij}\right)^2
\]
according to whether two entries have the same row, the same column, or neither, gives
\begin{equation}\label{eq:cross-identity}
 \sum_{\substack{i\ne k\\j\ne\ell}}x_{ij}x_{k\ell}
 =Q-R_2-C_2.
\end{equation}
Substituting \eqref{eq:cross-identity} into \eqref{eq:per-second} proves \eqref{eq:hessian}.  Since $0<c_n\le1$ and $Q>0$ for $X\ne0$, the right side is strictly negative.
\end{proof}

\begin{corollary}\label{cor:strict-local}
The matrix $U_n$ is a strict local maximizer of $\Phi$ on $K_n$.
\end{corollary}

\begin{proof}
For a tangent vector $X=(x_{ij})$, put
$u_i=\sum_jx_{ij}$ and $v_j=\sum_ix_{ij}$.  Since
$\sum_{i,j}x_{ij}=0$, one has
$\sum_i u_i=\sum_jv_j=0$.  Direct differentiation at $U_n$ gives
\[
 D\Phi(U_n)[X]
 =\sum_i u_i+\sum_jv_j
 -\frac{(n-1)!}{n^{n-1}}\sum_{i,j}x_{ij}=0.
\]
The tangent space of the affine hull of $K_n$ is finite-dimensional.  By
Lemma~\ref{lem:hessian}, the quadratic form given by the Hessian is uniformly
negative on its unit sphere.  The vanishing linear term and Taylor's formula
then give a relative neighborhood $\mathcal N$ of $U_n$ in $K_n$ such that
\[
 A\in\mathcal N,\ A\ne U_n
 \quad\Longrightarrow\quad
 \Phi(A)<\Phi(U_n).
\]
\end{proof}

\subsection{Averaging convergence and the interior classification}

\begin{lemma}[Convergence of cyclic pairwise averaging]\label{lem:averaging-convergence}
Let $c_1,\dots,c_n$ be vectors, and repeatedly replace pairs by their arithmetic mean according to a cyclic schedule containing every unordered pair.  Then all vectors converge to their common initial mean.
\end{lemma}

\begin{proof}
Let $\bar c=n^{-1}\sum_jc_j$, which is invariant, and define
\[
 V=\sum_j\|c_j-\bar c\|^2.
\]
Averaging $c_s,c_t$ decreases $V$ by
\begin{equation}\label{eq:variance-drop}
 \frac12\|c_s-c_t\|^2.
\end{equation}
Thus $V$ decreases to a limit $L\ge0$.

Assume $L>0$.  Take a convergent subsequence of the states at the beginnings of complete cycles.  The total decrease during each selected cycle tends to zero.  Since that decrease is a finite sum of the nonnegative quantities in \eqref{eq:variance-drop}, every pair difference at the moment it is averaged tends to zero.  Passing to the limiting cycle, every averaging operation makes no change.  Since every unordered pair occurs, all limiting vectors are equal, forcing $V=0$, a contradiction.  Hence $L=0$.
\end{proof}

\begin{theorem}[Positive-maximizer theorem]\label{thm:positive-maximizer}
For every $n\ge2$, every positive global maximizer of $\Phi$ on $K_n$ is $U_n$.
\end{theorem}

\begin{proof}
Let $A$ be a positive global maximizer.  Repeatedly average pairs of columns in a cyclic schedule containing every unordered pair.  By Lemma~\ref{lem:flat-average}, every matrix in the sequence is a positive global maximizer with the same value as $A$.  By Lemma~\ref{lem:averaging-convergence}, the sequence converges to a positive matrix $C$ with all columns equal.  Continuity shows that $C$ is also a global maximizer.

Now cyclically average pairs of rows of $C$.  Every matrix remains a positive global maximizer, and the sequence converges to $U_n$: the equal columns of $C$ each have sum one, and row averaging preserves every column sum.  Therefore
\begin{equation}\label{eq:U-global}
 \Phi(U_n)=\max_{K_n}\Phi.
\end{equation}

Let $Y_m$ be the row-averaging sequence converging to $U_n$.  Corollary~\ref{cor:strict-local} implies that $Y_m=U_n$ for all sufficiently large $m$, because every $Y_m$ has the value $\Phi(U_n)$.

We claim that one pair-averaging operation cannot map a distinct positive global maximizer $X$ to $U_n$.  If it did, the path \eqref{eq:pair-path} would have constant value by Lemma~\ref{lem:flat-average}, and points with $t\ne1/2$ arbitrarily close to $1/2$ would be distinct from $U_n$, arbitrarily close to it, and have value $\Phi(U_n)$.  This contradicts strict local maximality.  Hence the predecessor of $U_n$ under such an averaging operation is also $U_n$.

Applying this observation backward through the finite row-averaging segment gives $C=U_n$.  The column-averaging sequence converges to $U_n$ and has constant value, so strict local maximality again forces some finite column-averaging iterate to equal $U_n$.  Applying the same backward argument gives $A=U_n$.
\end{proof}

\begin{remark}
Theorem~\ref{thm:positive-maximizer} reproduces the positive-support conclusion historically due to Hwang, but its proof here is independent of \cite{Hwang1986}.
\end{remark}

\section{Joint deficits, subset sums, and Hall dilation}\label{sec:hall}

Let $A$ be a global maximizer of $\Phi$ on $K_n$.  Since $U_n$ is feasible,
\begin{equation}\label{eq:near-max}
 \Phi(A)\ge \Phi(U_n)=2-\gamma_n.
\end{equation}
Write
\begin{equation}\label{eq:deficits}
 \prod_i r_i=1-\rho,
 \qquad
 \prod_jc_j=1-\sigma,
 \qquad
 \per(A)=\gamma_n-\delta.
\end{equation}

\begin{lemma}[Joint deficit budget]\label{lem:joint-deficit}
The quantities in \eqref{eq:deficits} satisfy
\begin{equation}\label{eq:joint-budget}
 0\le\delta\le\gamma_n,
 \qquad
 \rho\ge0,
 \qquad
 \sigma\ge0,
 \qquad
 \rho+\sigma\le\delta.
\end{equation}
In particular, all row sums and column sums of $A$ are positive.
\end{lemma}

\begin{proof}
The row sums and column sums each have arithmetic mean one, so AM--GM gives
\[
 \prod_i r_i\le1,
 \qquad
 \prod_jc_j\le1.
\]
Thus $\rho,\sigma\ge0$.  Also
\[
 \per(A)
 =\prod_i r_i+\prod_jc_j-\Phi(A)
 \le2-(2-\gamma_n)=\gamma_n,
\]
which gives $0\le\delta\le\gamma_n$.  Substitution into \eqref{eq:near-max} yields
\[
 2-\rho-\sigma-\gamma_n+\delta\ge2-\gamma_n,
\]
and hence $\rho+\sigma\le\delta$.  If a row or column sum vanished, one of $\rho,\sigma$ would equal one, contradicting $\delta\le\gamma_n<1$.
\end{proof}

\subsection{Hall domination}

For a nonnegative matrix $S$ and index sets $I,J\subseteq[n]$, write
\[
 s_S(I,J)=\sum_{i\in I,\ j\in J}s_{ij}.
\]

\begin{lemma}[Hall domination]\label{lem:hall}
A nonnegative $n\times n$ matrix $S$ dominates a doubly stochastic matrix entrywise if and only if
\begin{equation}\label{eq:hall-condition}
 s_S(I,J)\ge |I|+|J|-n
\end{equation}
for every $I,J\subseteq[n]$ with $|I|+|J|>n$.

If equality holds in \eqref{eq:hall-condition} for some $I,J$, then every doubly stochastic $B\le S$ satisfies
\begin{equation}\label{eq:forced-block}
 B[I^c,J^c]=0.
\end{equation}
\end{lemma}

\begin{proof}
Construct a flow network with source-to-row capacities one, row-to-column capacities $s_{ij}$, and column-to-sink capacities one.  A flow of value $n$ is exactly a doubly stochastic matrix $B\le S$.  A cut having row set $I$ and column set $K$ on the source side has capacity
\[
 n-|I|+s_S(I,K^c)+|K|.
\]
By max-flow/min-cut, every cut has capacity at least $n$ exactly when \eqref{eq:hall-condition} holds, after writing $J=K^c$.

Let
\[
 p=|I|+|J|-n>0.
\]
Every doubly stochastic $B$ satisfies
\begin{equation}\label{eq:DS-cut-identity}
 s_B(I,J)-s_B(I^c,J^c)=p.
\end{equation}
If $s_S(I,J)=p$ and $B\le S$, then $s_B(I,J)\le p$.  Equation~\eqref{eq:DS-cut-identity} and nonnegativity give the reverse inequality, so equality holds and $s_B(I^c,J^c)=0$.
\end{proof}

\subsection{Relative-entropy subset estimates}

For $0<p,q<1$, let
\[
 D(p\Vert q)=p\log\frac pq+(1-p)\log\frac{1-p}{1-q}.
\]
We shall use the integral identity
\begin{equation}\label{eq:KL-integral}
 D(p\Vert q)=\int_p^q\frac{u-p}{u(1-u)}\,du.
\end{equation}
It follows by differentiating with respect to $q$.  In every subset-excess
application below one has $q\ge p$.  On that oriented interval the integrand
is nonnegative, and $u(1-u)\le1/4$ gives
\[
 D(p\Vert q)\ge4\int_p^q(u-p)\,du=2(q-p)^2.
\]
The case $q<p$ follows by the same argument after reversing the limits (or
by symmetry), so the binary Pinsker bound is
\begin{equation}\label{eq:pinsker}
 D(p\Vert q)\ge2(p-q)^2.
\end{equation}

For $7\le n\le14$, define $\eta_n$ by
\begin{equation}\label{eq:eta-table}
\begin{array}{c|cccccccc}
 n&7&8&9&10&11&12&13&14\\ \hline
 \eta_n&1/47&1/80&1/135&1/230&1/350&1/650&1/1100&1/1800.
\end{array}
\end{equation}
For $0\le k\le n-1$, define
\[
 \alpha_{n,0}=0
\]
and, for $k\ge1$,
\begin{equation}\label{eq:alpha-mid}
 \alpha_{n,k}=
 \begin{cases}
  2n\left(\dfrac{k}{n}+\eta_n\right)
     \left(1-\dfrac{k}{n}-\eta_n\right),&2k<n,\\[3mm]
  \dfrac{2k(n-k)}n,&2k\ge n.
 \end{cases}
\end{equation}
The exact arithmetic in Appendix~\ref{app:certificates} verifies
\begin{equation}\label{eq:eta-localization}
 \eta_n^2>\frac{\gamma_n}{2n(1-\gamma_n)}
\end{equation}
and
\begin{equation}\label{eq:eta-half}
 \frac{\lfloor(n-1)/2\rfloor}{n}+\eta_n<\frac12
\end{equation}
for $7\le n\le14$.

\begin{lemma}[Subset excess, dimensions $7$ through $14$]\label{lem:subset-mid}
Let $7\le n\le14$.  If $E$ is a set of $k$ rows, $1\le k\le n-1$, and
\[
 \sum_{i\in E}r_i=k+\eps,
 \qquad \eps>0,
\]
then
\begin{equation}\label{eq:subset-mid}
 \eps^2\le\frac{\alpha_{n,k}\rho}{1-\delta}.
\end{equation}
The analogous estimate holds for column subsets, with $\sigma$ in place of $\rho$.
\end{lemma}

\begin{proof}
AM--GM applied separately on $E$ and $E^c$ gives
\begin{equation}\label{eq:subset-AMGM}
 1-\rho
 \le
 \left(1+\frac{\eps}{k}\right)^k
 \left(1-\frac{\eps}{n-k}\right)^{n-k}.
\end{equation}
Put
\[
 p=\frac{k}{n},
 \qquad
 q=\frac{k+\eps}{n}.
\]
Taking logarithms in \eqref{eq:subset-AMGM},
\begin{equation}\label{eq:KL-upper}
 nD(p\Vert q)\le-\log(1-\rho)
 \le\frac{\rho}{1-\rho}
 \le\frac{\rho}{1-\delta}.
\end{equation}
Pinsker's bound and Lemma~\ref{lem:joint-deficit} imply
\begin{equation}\label{eq:q-localized}
 (q-p)^2
 \le\frac{\gamma_n}{2n(1-\gamma_n)}
 <\eta_n^2.
\end{equation}

If $p<1/2$, then \eqref{eq:eta-half} and \eqref{eq:q-localized} imply $q<p+\eta_n<1/2$.  Hence
\[
 u(1-u)\le(p+\eta_n)(1-p-\eta_n)
 \qquad(p\le u\le q).
\]
Using \eqref{eq:KL-integral},
\[
 nD(p\Vert q)
 \ge\frac{\eps^2}
 {2n(p+\eta_n)(1-p-\eta_n)}
 =\frac{\eps^2}{\alpha_{n,k}}.
\]
If $p\ge1/2$, then $u(1-u)\le p(1-p)$ for $u\in[p,q]$, giving the same conclusion with $\alpha_{n,k}=2np(1-p)$.  Combine this lower bound with \eqref{eq:KL-upper}.
\end{proof}

For $n=15$, a sharper singleton constant is needed.  Set
\begin{equation}\label{eq:alpha15}
 \alpha_{15,0}=0,
 \qquad
 \alpha_{15,1}=\frac{195}{98},
 \qquad
 \alpha_{15,k}=\frac{15}{2}\quad(2\le k\le14).
\end{equation}

\begin{lemma}[Subset excess in dimension $15$]\label{lem:subset15}
The conclusion of Lemma~\ref{lem:subset-mid} holds for $n=15$ with the constants in \eqref{eq:alpha15}.
\end{lemma}

\begin{proof}
The bound $\alpha_{15,k}=15/2$ follows from \eqref{eq:pinsker} exactly as in the preceding proof.

Suppose $k=1$.  The exact comparison $\gamma_{15}<1/300000$, verified in Appendix~\ref{app:certificates}, gives
\[
 \eps^2
 \le\frac{15\gamma_{15}}{2(1-\gamma_{15})}
 <\frac{15}{2\cdot299999}<\frac1{196}.
\]
Thus $0<\eps<1/14$, and with $p=1/15$ and $q=(1+\eps)/15$,
\[
 \frac1{15}<q<\frac1{14}.
\]
On $[1/15,q]$,
\[
 u(1-u)\le\frac{13}{196}.
\]
Equation~\eqref{eq:KL-integral} therefore gives
\[
 15D\left(\frac1{15}\middle\Vert q\right)
 \ge\frac{98}{195}\eps^2.
\]
Combining this with \eqref{eq:KL-upper} proves
\[
 \eps^2\le\frac{195}{98}\frac{\rho}{1-\delta}.
\]
\end{proof}

\subsection{The critical Hall deficit}

Define
\begin{equation}\label{eq:tau}
 \tau=
 \max\left\{
 0,
 \max_{\substack{I,J\subseteq[n]\\|I|+|J|>n}}
 \left(
 1-\frac{s_A(I,J)}{|I|+|J|-n}
 \right)
 \right\}.
\end{equation}
Assume $\tau>0$, and choose a critical pair $I,J$.  Put
\begin{equation}\label{eq:abp}
 a=|I^c|,
 \qquad
 b=|J^c|,
 \qquad
 p=n-a-b\ge1.
\end{equation}

\begin{lemma}[Hall-deficit estimate]\label{lem:hall-deficit}
With the constants $\alpha_{n,k}$ defined above,
\begin{equation}\label{eq:tau-bound}
 \tau^2\le c_{n,a,b}\frac{\delta}{1-\delta},
 \qquad
 c_{n,a,b}=
 \frac{\alpha_{n,a}+\alpha_{n,b}}{p^2}.
\end{equation}
Whenever the right side is less than one, the matrix
\begin{equation}\label{eq:S-dilation}
 S=\frac{A}{1-\tau}
\end{equation}
dominates a doubly stochastic matrix $B$, and every such $B$ satisfies
\begin{equation}\label{eq:B-block}
 B[I^c,J^c]=0.
\end{equation}
Every zero entry of $A$ is also a zero entry of $B$.
\end{lemma}

\begin{proof}
Let
\[
 e=\sum_{i\in I^c}r_i-a,
 \qquad
 f=\sum_{j\in J^c}c_j-b,
 \qquad
 x=s_A(I^c,J^c)\ge0.
\]
Inclusion--exclusion gives
\[
 s_A(I,J)=p-e-f+x.
\]
Criticality therefore implies
\begin{equation}\label{eq:critical-excess}
 \tau p=e+f-x\le e_++f_+.
\end{equation}
By the subset estimates,
\[
 e_+^2\le\frac{\alpha_{n,a}\rho}{1-\delta},
 \qquad
 f_+^2\le\frac{\alpha_{n,b}\sigma}{1-\delta}.
\]
Cauchy--Schwarz and \eqref{eq:joint-budget} give
\[
 (e_++f_+)^2
 \le\frac{(\alpha_{n,a}+\alpha_{n,b})(\rho+\sigma)}{1-\delta}
 \le\frac{(\alpha_{n,a}+\alpha_{n,b})\delta}{1-\delta}.
\]
Together with \eqref{eq:critical-excess}, this proves \eqref{eq:tau-bound}.

By the definition of $\tau$, the matrix $S$ satisfies every Hall inequality \eqref{eq:hall-condition}.  Lemma~\ref{lem:hall} gives a doubly stochastic $B\le S$.  Equality holds for the critical pair, so \eqref{eq:B-block} follows.  Finally, $S$ has every zero of $A$, and $B\le S$.
\end{proof}

If $\tau=0$, Lemma~\ref{lem:hall} gives a doubly stochastic $B\le A$.  Since both matrices have total sum $n$, necessarily $B=A$.  Thus a boundary maximizer with $\tau=0$ would itself be doubly stochastic with a zero, and Corollary~\ref{cor:one-zero} would give
\[
 \per(A)\ge\kappa_n>\gamma_n,
\]
contrary to Lemma~\ref{lem:joint-deficit}.  Hence every boundary maximizer must have $\tau>0$.

\subsection{A common normalized block-variation identity}
The stationarity identities used in Sections~\ref{sec:mid-high},
\ref{sec:low}, \ref{sec:n7}, and \ref{sec:n6} are instances of the same
calculation.  Let $E\subseteq[n]\times[n]$ be a set of entries of total mass
$m_E$, let
\[
 e_i=\sum_{j:(i,j)\in E}a_{ij},\qquad
 f_j=\sum_{i:(i,j)\in E}a_{ij},
\]
and, for a permutation $\sigma$, let
$k_E(\sigma)=|\{i:(i,\sigma(i))\in E\}|$.  Multiply the entries in $E$ by
$t$ and then multiply the whole matrix by
\[
 \lambda(t)=\frac{n}{n+(t-1)m_E}
\]
so that the total mass remains $n$.  At a global maximizer, differentiation
at $t=1$ gives
\begin{equation}\label{eq:normalized-block-variation}
 0=
 R\sum_i\frac{e_i}{r_i}
 +C\sum_j\frac{f_j}{c_j}
 -\sum_{\sigma\in S_n}k_E(\sigma)
    \prod_i a_{i,\sigma(i)}
 -m_E\Phi(A),
\end{equation}
where $R=\prod_i r_i$ and $C=\prod_jc_j$.  Indeed, the first three terms
are the derivative before normalization, while homogeneity of degree $n$
and $n\lambda'(1)=-m_E$ give the last term.  Each later stationarity lemma
specifies $E$, its row and column masses, and the matching count
$k_E(\sigma)$; thus no normalization factor is implicit.

\subsection{Scalar exclusion lemmas}

\begin{lemma}[First-order scalar criterion]\label{lem:first-order}
Suppose
\[
 \per(B)\ge L,
 \qquad
 \tau^2\le c\frac{\delta}{1-\delta},
 \qquad
 0\le\delta\le\gamma_n,
 \qquad
 0\le\tau<1.
\]
Then
\begin{equation}\label{eq:first-margin}
 L(1-\tau)^n-(\gamma_n-\delta)
 \ge
 M_n^{(1)}(L,c),
\end{equation}
where
\begin{equation}\label{eq:M1}
 M_n^{(1)}(L,c)
 =L-\gamma_n-
 \frac{n^2cL^2}{4(1-\gamma_n)}.
\end{equation}
In particular, $M_n^{(1)}(L,c)>0$ contradicts
\[
 \per(A)=\gamma_n-\delta
 \ge L(1-\tau)^n.
\]
\end{lemma}

\begin{proof}
Bernoulli's inequality gives $(1-\tau)^n\ge1-n\tau$.  Hence
\[
 \begin{aligned}
 L(1-\tau)^n-(\gamma_n-\delta)
 &\ge L-\gamma_n+\delta-nL\tau\\
 &\ge L-\gamma_n+\delta
 -\frac{nL\sqrt c}{\sqrt{1-\gamma_n}}\sqrt\delta.
 \end{aligned}
\]
Complete the square in $\sqrt\delta$.
\end{proof}

\begin{lemma}[Second-order scalar criterion]\label{lem:second-order}
Suppose, in addition, that $0\le\tau\le T$ and
\begin{equation}\label{eq:quadratic-binomial}
 (1-t)^n\ge1-nt+\lambda t^2
 \qquad(0\le t\le T).
\end{equation}
Then
\begin{equation}\label{eq:second-margin}
 L(1-\tau)^n-(\gamma_n-\delta)
 \ge
 M_n^{(2)}(L,c),
\end{equation}
where
\begin{equation}\label{eq:M2}
 M_n^{(2)}(L,c)
 =L-\gamma_n-
 \frac{n^2L^2}
 {4\left(\lambda L+(1-\gamma_n)/c\right)}.
\end{equation}
\end{lemma}

\begin{proof}
The Hall bound implies
\[
 \delta\ge\frac{1-\gamma_n}{c}\tau^2.
\]
Therefore
\[
 \begin{aligned}
 L(1-\tau)^n-(\gamma_n-\delta)
 &\ge L-\gamma_n-nL\tau
 +\left(\lambda L+\frac{1-\gamma_n}{c}\right)\tau^2.
 \end{aligned}
\]
The minimum of the quadratic on the right is \eqref{eq:M2}.
\end{proof}

\section{Boundary exclusion in dimensions eleven through fifteen}\label{sec:mid-high}

We now apply the common framework.  Throughout this section, $A$ denotes a hypothetical boundary global maximizer.  As observed after Lemma~\ref{lem:hall-deficit}, necessarily $\tau>0$.

\subsection{Dimensions eleven through fourteen: ordinary Hall pairs}

For $11\le n\le14$, let
\[
 \mathcal P_n=
 \{(a,b):0\le a\le b,\ a+b\le n-1\}.
\]
Transposition permits the restriction $a\le b$.  For $(a,b)\in\mathcal P_n$, put
\begin{equation}\label{eq:mid-c}
 p=n-a-b,
 \qquad
 c_{n,a,b}=
 \frac{\alpha_{n,a}+\alpha_{n,b}}{p^2},
\end{equation}
and
\begin{equation}\label{eq:mid-L}
 \Lambda_{n,a,b}=
 \begin{cases}
  \kappa_n,&a=0,\\[1mm]
  v_{n-a}G(n-b)^a,&a\ge1.
 \end{cases}
\end{equation}
If $a=0$, the forced Hall block is empty, but $B$ inherits a zero from $A$, so Corollary~\ref{cor:one-zero} gives the first line.  If $a\ge1$, Corollary~\ref{cor:weak-block} gives the second.

The constants in \eqref{eq:alpha-mid} satisfy $\alpha_{n,k}\le n/2$: in the first branch this is $2n x(1-x)\le n/2$, and in the second it is $2k(n-k)/n\le n/2$.  Hence $c_{n,a,b}\le n$.  Moreover
\[
 \frac{\gamma_{n+1}}{\gamma_n}=\left(\frac{n}{n+1}\right)^n<1,
\]
so $\gamma_n\le\gamma_{11}<1/1000$ for $11\le n\le14$.  Therefore
\[
 \tau^2\le\frac{n\gamma_n}{1-\gamma_n}
 <\frac{14}{999}<1.
\]
Thus the Hall dilation is valid.

The following proposition is a finite exact-rational calculation.  Its formulas are explicit in \eqref{eq:mid-c}, \eqref{eq:mid-L}, and \eqref{eq:M1}; Appendix~\ref{app:code} gives a complete standard-library implementation using rational arithmetic only.

\begin{proposition}[Exact ordinary margins]\label{prop:mid-margins}
For the indicated sets of pairs,
\[
 M_n^{(1)}(\Lambda_{n,a,b},c_{n,a,b})>0.
\]
More precisely, the minimum values are as follows:
\begin{equation}\label{eq:mid-margin-table}
\begin{array}{c|c|c}
 n&\text{minimizing pair}&\text{minimum (decimal for scale)}\\ \hline
 11&(0,9)&1.919523\times10^{-7}\\
 12&(0,11)&2.454520\times10^{-8}\\
 13&(0,12)&3.656538\times10^{-8}\\
 14&(0,13)&1.708143\times10^{-8}.
\end{array}
\end{equation}
The decimal values in this and later summary tables are only readability
aids.  The proof compares the exact fractions recorded by the checker and in
\texttt{verification\_manifest.json}.
For $n=11$, the pair $(0,10)$ is excluded from this proposition and is handled separately below.
\end{proposition}

\begin{proof}
Substitute the rational definitions into \eqref{eq:M1} for every pair in $\mathcal P_n$, omitting $(0,10)$ when $n=11$.  There are respectively $35,42,49,56$ cases.  Clearing denominators gives positive integers in every case.  The exact enumeration in Appendix~\ref{app:code} also compares the resulting fractions and returns the minima displayed in \eqref{eq:mid-margin-table}.  No rounded quantity is used.
\end{proof}

\begin{corollary}\label{cor:ordinary-excluded}
All critical Hall pairs in dimensions $12,13,14$ are impossible.  In dimension $11$, every critical pair except $(a,b)=(0,10)$ and its transpose is impossible.
\end{corollary}

\begin{proof}
For a critical pair, Lemma~\ref{lem:hall-deficit} gives $B$ and the bound \eqref{eq:tau-bound}; the appropriate permanent estimate is \eqref{eq:mid-L}.  Proposition~\ref{prop:mid-margins} and Lemma~\ref{lem:first-order} contradict \eqref{eq:per-comparison}, where
\begin{equation}\label{eq:per-comparison}
 \per(A)=(1-\tau)^n\per(S)
 \ge(1-\tau)^n\per(B).
\end{equation}
\end{proof}

\subsection{The exceptional deficient column in dimension eleven}

Assume $n=11$ and $(a,b)=(0,10)$.  Then $I=[11]$, while $J$ consists of one column, say column $j$, and
\begin{equation}\label{eq:deficient-column}
 c_j=1-\tau.
\end{equation}
The complementary set of ten columns has total sum $10+\tau$.  Lemma~\ref{lem:subset-mid} gives
\begin{equation}\label{eq:n11-column}
 \tau^2\le\frac{\alpha_{11,10}\sigma}{1-\delta}.
\end{equation}
We obtain a second use of $\tau$ from stationarity.

\begin{lemma}[Deficient-column stationarity]\label{lem:n11-stationarity}
Let $A$ be a global maximizer in dimension $11$.  Suppose its row sums are positive and a column has sum $1-\tau<1$.  If
\[
 d=\max_i(r_i-1)_+,
\]
then
\begin{equation}\label{eq:n11-d}
 d\ge(1-\gamma_{11})\tau.
\end{equation}
\end{lemma}

\begin{proof}
Apply the normalized block-variation identity
\eqref{eq:normalized-block-variation} to the deficient column.  Here the
block mass is $m_E=1-\tau$, its row masses are $a_{ij}$, its only nonzero
column mass is $1-\tau$, and every permanent monomial uses exactly one
entry of that column.  With
\[
 u=\sum_i\frac{a_{ij}}{r_i},
 \qquad
 K=\frac{\Phi(A)}{\prod_i r_i},
\]
identity \eqref{eq:normalized-block-variation} becomes
\[
 \left(\prod_i r_i\right)u
 +\prod_jc_j-\per(A)
 -(1-\tau)\Phi(A)=0.
\]
Since $\prod_jc_j-\per(A)=\Phi(A)-\prod_i r_i$, this simplifies to
\begin{equation}\label{eq:n11-u}
 u=1-K\tau.
\end{equation}
Now put $w_i=a_{ij}/r_i$.  Then $w_i\ge0$, $\sum_iw_i=u$, and
\[
 (1-\tau)-u
 =\sum_iw_i(r_i-1)
 =(K-1)\tau.
\]
The right side is positive.  Moreover $0<u<1$ by \eqref{eq:n11-u}, so
\[
 (K-1)\tau
 \le d\sum_iw_i
 =du\le d.
\]
Finally, $\Phi(A)\ge2-\gamma_{11}$ and $\prod_i r_i\le1$, so $K\ge2-\gamma_{11}$.
\end{proof}

Applying Lemma~\ref{lem:subset-mid} to a singleton row attaining $d$ and using Lemma~\ref{lem:n11-stationarity},
\begin{equation}\label{eq:n11-rho}
 \rho\ge
 (1-\delta)\frac{(1-\gamma_{11})^2}{\alpha_{11,1}}\tau^2.
\end{equation}
Equation~\eqref{eq:n11-column} gives
\begin{equation}\label{eq:n11-sigma}
 \sigma\ge
 (1-\delta)\frac1{\alpha_{11,10}}\tau^2.
\end{equation}
Together with $\rho+\sigma\le\delta$,
\begin{equation}\label{eq:n11-cstar-bound}
 \tau^2\le c_*\frac{\delta}{1-\delta},
\end{equation}
where
\begin{equation}\label{eq:n11-cstar}
 c_*=
 \left(
 \frac{(1-\gamma_{11})^2}{\alpha_{11,1}}
 +\frac1{\alpha_{11,10}}
 \right)^{-1}
 =\frac{1540632960668022501507138780}
 {1671236290061169784540151329}.
\end{equation}
Exact subtraction gives
\begin{equation}\label{eq:n11-cstar-gap}
 \frac{93}{100}-c_*
 =\frac{1361678908886539811520195597}
 {167123629006116978454015132900}>0.
\end{equation}
Since $B$ inherits a zero, $\per(B)\ge\kappa_{11}$.  For fixed $L>0$,
formula~\eqref{eq:M1} shows that $M_n^{(1)}(L,c)$ is strictly decreasing in
$c$.  Hence the certified upper bound $c_*<93/100$ permits the larger,
rationally convenient value $93/100$.  Finally,
\begin{equation}\label{eq:n11-special-margin}
 M_{11}^{(1)}\left(\kappa_{11},\frac{93}{100}\right)>0,
 \qquad
 M_{11}^{(1)}\left(\kappa_{11},\frac{93}{100}\right)
 \approx1.250887\times10^{-7}.
\end{equation}
The exact numerator and denominator are recorded literally in the
machine-readable verification manifest and checked by the standard-library
verifier.
Lemma~\ref{lem:first-order} gives a contradiction.  Therefore dimension $11$ has no boundary global maximizer.

\subsection{Dimension fifteen: a thin--thick dichotomy}

Let $n=15$, and retain the notation $(a,b,p)$ from \eqref{eq:abp}.  Call a critical pair \emph{thin} if
\[
 a=0,
 \quad b=0,
 \quad\text{or}\quad
 \min\{a,b\}=1.
\]
Otherwise call it \emph{thick}.  Since $\alpha_{15,k}\le15/2$ for every $k$, Lemma~\ref{lem:hall-deficit} and \eqref{eq:n15-gamma-bounds} give
\[
 \tau^2\le\frac{15\gamma_{15}}{1-\gamma_{15}}
 <\frac{15}{299999}<1.
\]
Thus the Hall dilation and the doubly stochastic matrix $B$ are available in both cases.

For a thin pair, \eqref{eq:alpha15} gives
\[
 \alpha_{15,a}+\alpha_{15,b}
 \le\frac{195}{98}+\frac{15}{2}
 =\frac{465}{49}.
\]
Since $p\ge1$,
\begin{equation}\label{eq:n15-thin-tau}
 \tau^2\le\frac{465}{49}\frac{\delta}{1-\delta}.
\end{equation}
The matrix $B$ inherits a zero, so
\begin{equation}\label{eq:mu1}
 \per(B)\ge\mu_1,
 \qquad
 \mu_1:=v_{14}G(14)
 =\frac{13!\,13^{13}}{14^{26}}.
\end{equation}

For a thick pair, $a,b\ge2$, so $B$ contains a $2\times2$ zero block.  Thus
\begin{equation}\label{eq:mu2}
 \per(B)\ge\mu_2,
 \qquad
 \mu_2:=v_{13}G(13)^2
 =\frac{13!\,12^{24}}{13^{37}}.
\end{equation}
Also
\begin{equation}\label{eq:n15-thick-tau}
 \tau^2\le15\frac{\delta}{1-\delta}.
\end{equation}
The exact comparisons
\begin{equation}\label{eq:n15-gamma-bounds}
 \gamma_{15}<\frac1{300000},
 \qquad
 \gamma_{15}<\frac{29863}{10^{10}},
\end{equation}
\begin{equation}\label{eq:n15-mu1-bounds}
 \frac{29937}{10^{10}}<\mu_1<\frac3{10^6},
\end{equation}
and
\begin{equation}\label{eq:n15-mu2-bounds}
 \frac{301}{10^8}<\mu_2<\frac4{10^6}
\end{equation}
are certified by the six positive integers listed in Appendix~\ref{app:certificates}.

For the thin case,
\[
 \frac{225}{4}\frac{465}{49}
 =\frac{104625}{196}<534.
\]
Equations \eqref{eq:n15-gamma-bounds} and \eqref{eq:n15-mu1-bounds} give
\[
 \mu_1-\gamma_{15}>\frac{74}{10^{10}}
\]
and
\[
 \frac{225(465/49)\mu_1^2}{4(1-\gamma_{15})}
 <534\cdot\frac9{10^{12}}\cdot\frac{300000}{299999}
 <\frac5{10^9}.
\]
The final strict inequality follows after clearing denominators; its integer gap is
\[
 58\,195\,000\,000\,000\,000>0.
\]
Therefore
\begin{equation}\label{eq:n15-thin-margin}
 M_{15}^{(1)}\left(\mu_1,\frac{465}{49}\right)
 >\frac{74}{10^{10}}-\frac5{10^9}
 =\frac{24}{10^{10}}>0.
\end{equation}

For the thick case, $\gamma_{15}<299/10^8$, so
\[
 \mu_2-\gamma_{15}>\frac2{10^8}.
\]
Also
\[
 \frac{225\cdot15\mu_2^2}{4(1-\gamma_{15})}
 <\frac{3375}{4}\cdot\frac{16}{10^{12}}
 \cdot\frac{300000}{299999}
 <\frac{14}{10^9},
\]
where the final cleared-denominator gap is
\[
 149\,986\,000\,000\,000\,000>0.
\]
Hence
\begin{equation}\label{eq:n15-thick-margin}
 M_{15}^{(1)}(\mu_2,15)
 >\frac2{10^8}-\frac{14}{10^9}
 =\frac6{10^9}>0.
\end{equation}
Both cases contradict Lemma~\ref{lem:first-order}.  Thus dimension $15$ has no boundary global maximizer.

\begin{theorem}[Boundary exclusion for $11\le n\le15$]\label{thm:boundary-mid-high}
For every $n\in\{11,12,13,14,15\}$, no global maximizer of $\Phi$ on $K_n$ has a zero entry.
\end{theorem}

\begin{proof}
If $\tau=0$, the contradiction following Lemma~\ref{lem:hall-deficit} applies.  If $\tau>0$, Corollary~\ref{cor:ordinary-excluded} handles dimensions $12,13,14$ and all ordinary pairs in dimension $11$; equations \eqref{eq:n11-cstar-bound}--\eqref{eq:n11-special-margin} handle the exceptional pair.  The thin--thick argument handles dimension $15$.
\end{proof}

\section{Boundary exclusion in dimensions eight, nine, and ten}\label{sec:low}

For $n\in\{8,9,10\}$, define
\begin{equation}\label{eq:low-parameters}
\begin{array}{c|ccc}
 n&8&9&10\\ \hline
 T_n&1/7&1/10&1/10\\
 \lambda_n&20&27&33.
\end{array}
\end{equation}
As before, let $A$ be a hypothetical boundary global maximizer.  The constants $\alpha_{n,k}$ satisfy $\alpha_{n,k}\le n/2$, so Lemma~\ref{lem:hall-deficit} gives
\begin{equation}\label{eq:low-tau-range}
 \tau^2\le\frac{n\gamma_n}{1-\gamma_n}<T_n^2.
\end{equation}
The exact positive gaps in \eqref{eq:low-tau-range} are
\begin{equation}\label{eq:tau-range-gaps}
\begin{array}{c|c}
 n&T_n^2-\dfrac{n\gamma_n}{1-\gamma_n}\\ \hline
 8&\dfrac{7277}{6407093}\\[2mm]
 9&\dfrac{746489}{477848900}\\[2mm]
 10&\dfrac{994933}{156193300}.
\end{array}
\end{equation}

\begin{lemma}[Quadratic binomial bounds]\label{lem:binomial-low}
For $n\in\{8,9,10\}$ and $0\le t\le T_n$,
\begin{equation}\label{eq:low-binomial}
 (1-t)^n\ge1-nt+\lambda_nt^2.
\end{equation}
\end{lemma}

\begin{proof}
From the binomial expansion,
\[
 (1-t)^n
 \ge1-nt+\binom n2t^2-\binom n3t^3.
\]
Indeed, after the cubic term the alternating terms have decreasing absolute values on the stated intervals: for $k\ge3$,
\[
 \frac{\binom n{k+1}t^{k+1}}{\binom nk t^k}
 =\frac{n-k}{k+1}t<1.
\]
Thus the remaining tail starts positive and is nonnegative.  Finally,
\[
 \binom82-\binom83\frac17=20,
\]
\[
 \binom92-\binom93\frac1{10}=\frac{138}{5}>27,
\]
and
\[
 \binom{10}2-\binom{10}3\frac1{10}=33.
\]
\end{proof}

\subsection{Two-sided critical pairs}

Suppose $a,b\ge1$.  By transposition, assume $a\le b$.  Put
\[
 p=n-a-b,
 \qquad
 c_{n,a,b}=
 \frac{\alpha_{n,a}+\alpha_{n,b}}{p^2},
\]
and use the sharp block bound
\[
 L^{\sharp}_{n,a,b}=\frac{v_{n-a}v_{n-b}}{v_p}.
\]
Lemma~\ref{lem:second-order} applies with $\lambda=\lambda_n$.

\begin{proposition}[Exact two-sided margins]\label{prop:low-two-sided}
For every
\[
 n\in\{8,9,10\},
 \qquad
 1\le a\le b,
 \qquad
 a+b\le n-1,
\]
one has
\[
 M_n^{(2)}(L^{\sharp}_{n,a,b},c_{n,a,b})>0.
\]
The minimum occurs at $(a,b)=(1,1)$ and equals
\begin{equation}\label{eq:low-two-sided-table}
\begin{array}{c|c}
 n&\text{minimum (decimal for scale)}\\ \hline
 8&1.379790\times10^{-5}\\
 9&5.710428\times10^{-6}\\
 10&1.982111\times10^{-6}.
\end{array}
\end{equation}
\end{proposition}

\begin{proof}
There are respectively $12,16,20$ unordered pairs.  Substitute the exact rational definitions into \eqref{eq:M2}, clear denominators, and compare the resulting positive fractions.  Appendix~\ref{app:code} performs exactly this finite calculation.
\end{proof}

Thus no two-sided critical pair can occur.

\subsection{One-sided cuts and stationarity}

It remains to consider a one-sided critical pair.  After transposition,
assume $a=0$.  The case $b=0$ would give $I=J=[n]$ and Hall deficit zero,
so a critical pair with $\tau>0$ necessarily has
$1\le b\le n-1$.  Thus $I=[n]$ and $J$ is a set of
\[
 p=n-b
\]
columns satisfying
\begin{equation}\label{eq:one-sided-mass}
 \sum_{j\in J}c_j=p(1-\tau).
\end{equation}
The corresponding columns of $S=A/(1-\tau)$ have total sum $p$.  Since $B\le S$ is doubly stochastic, those same $p$ columns of $B$ also have total sum $p$.  Hence entrywise saturation holds:
\begin{equation}\label{eq:saturation}
 B[:,J]=S[:,J].
\end{equation}
Define
\begin{equation}\label{eq:beta}
 m_i=\sum_{j\in J}b_{ij},
 \qquad
 \beta_i=\frac{m_i}{p}.
\end{equation}
Then $\beta=(\beta_i)$ is a probability vector.

\begin{lemma}[One-sided stationarity identity]\label{lem:one-sided-stationarity}
Put
\[
 K=\frac{\Phi(A)}{\prod_i r_i}
\]
and
\[
 D=1-\sum_i\frac{\beta_i}{r_i}.
\]
Then
\begin{equation}\label{eq:D-lower}
 D=\frac{(K-1)\tau}{1-\tau}
 \ge\frac{(1-\gamma_n)\tau}{1-\tau}>0.
\end{equation}
\end{lemma}

\begin{proof}
Apply \eqref{eq:normalized-block-variation} to the $p$ columns in $J$.
Their total mass is $m_E=p(1-\tau)$; the row masses are
$\sum_{j\in J}a_{ij}$; each selected column has full column mass; and every
permanent monomial uses exactly $p$ selected entries.  With
\[
 u=\sum_i\frac{\sum_{j\in J}a_{ij}}{r_i},
\]
the general identity becomes
\[
 \left(\prod_i r_i\right)u
 +p\left(\prod_jc_j-\per(A)\right)
 -p(1-\tau)\Phi(A)=0.
\]
Using $\prod_jc_j-\per(A)=\Phi(A)-\prod_i r_i$,
\begin{equation}\label{eq:one-sided-u}
 u=p-pK\tau.
\end{equation}
By saturation,
\[
 \sum_{j\in J}a_{ij}=(1-\tau)m_i.
\]
Thus \eqref{eq:one-sided-u}, divided by $p$, becomes
\[
 (1-\tau)\sum_i\frac{\beta_i}{r_i}=1-K\tau,
\]
which is the asserted identity.  Finally,
\[
 K\ge\Phi(A)\ge2-\gamma_n,
\]
because $\prod_i r_i\le1$.
\end{proof}

\subsection{Row-product energy}

Define
\begin{equation}\label{eq:d0}
 d_0^2=\frac{n\gamma_n}{2(1-\gamma_n)}.
\end{equation}
For $n=8,9,10$, exact arithmetic gives $d_0<1/10$; the positive gaps are
\begin{equation}\label{eq:d0-gaps}
\begin{array}{c|c}
 n&\dfrac1{100}-d_0^2\\ \hline
 8&\dfrac{4757}{13075700}\\[2mm]
 9&\dfrac{2762489}{477848900}\\[2mm]
 10&\dfrac{1278433}{156193300}.
\end{array}
\end{equation}

\begin{lemma}[Row-product energy]\label{lem:row-energy}
For $n\in\{8,9,10\}$, put
\[
 y_i=1-\frac1{r_i},
 \qquad
 C_n=\frac{(1-\gamma_n)(1-d_0)^2}{2}.
\]
Then
\begin{equation}\label{eq:row-energy}
 \rho\ge C_n\sum_i y_i^2,
 \qquad C_n>\frac25.
\end{equation}
In particular, $\rho>\frac25\sum_i y_i^2$ whenever
$(y_1,\dots,y_n)\ne0$.
\end{lemma}

\begin{proof}
First, every row sum satisfies $r_i\ge1-d_0$.  Indeed, if $r_i=1-d<1$, the complementary $n-1$ rows have positive excess $d$.  Lemma~\ref{lem:subset-mid} and $\alpha_{n,n-1}\le n/2$ give
\[
 d^2\le\frac{n\rho}{2(1-\delta)}
 \le\frac{n\gamma_n}{2(1-\gamma_n)}=d_0^2.
\]

For every $r\ge1-d_0$,
\begin{equation}\label{eq:log-energy}
 r-1-\log r
 \ge\frac{(1-d_0)^2}{2}
 \left(1-\frac1r\right)^2.
\end{equation}
To see this, set $y=1-1/r$.  Then
\[
 r-1-\log r=\int_0^y\frac{t}{(1-t)^2}\,dt.
\]
For $y\ge0$, the integrand is at least $t$; for $y<0$, reverse the integral and use $(1-t)^{-2}\ge(1-d_0)^2$ on $[y,0]$.

Since $\sum_i(r_i-1)=0$,
\[
 -\log(1-\rho)
 =\sum_i(r_i-1-\log r_i)
 \ge\frac{(1-d_0)^2}{2}\sum_i y_i^2.
\]
Also
\[
 \rho=1-(1-\rho)
 \ge(1-\rho)(-\log(1-\rho))
 \ge(1-\gamma_n)(-\log(1-\rho)).
\]
For these dimensions, $\gamma_n<1/81$ and $d_0<1/10$, so
\[
 C_n=\frac{(1-\gamma_n)(1-d_0)^2}{2}
 >\frac{80}{81}\cdot\frac{81}{100}\cdot\frac12
 =\frac25.
\]
This proves the non-strict $C_n$ bound.  If $y\ne0$, then
$\sum_i y_i^2>0$ and the strict comparison with $2/5$ follows.
\end{proof}

Let
\begin{equation}\label{eq:beta-s}
 s=\sum_i\left(\beta_i-\frac1n\right)^2.
\end{equation}
By the harmonic-mean inequality,
\[
 \sum_i\frac1{r_i}\ge\frac{n^2}{\sum_i r_i}=n,
\]
so $\sum_i y_i\le0$.  Therefore
\[
 D=\sum_i\beta_i y_i
 \le
 \sum_i\left(\beta_i-\frac1n\right)y_i
 \le\sqrt{s}\left(\sum_i y_i^2\right)^{1/2}.
\]
Since $D>0$, one has $s>0$ and $y\ne0$.  Thus the strict form of
Lemma~\ref{lem:row-energy} applies, and \eqref{eq:D-lower} implies
\begin{equation}\label{eq:rho-D-s}
 \rho>\frac25\frac{D^2}{s}
 \ge
 \frac25\frac{(1-\gamma_n)^2\tau^2}
 {(1-\tau)^2s}.
\end{equation}

The $p$ columns in $J$ are probability vectors whose average is $\beta$.  Convexity of squared Euclidean distance shows that at least one of them, say $b^{(j)}$, satisfies
\[
 s(b^{(j)})\ge s.
\]
Theorem~\ref{thm:column-entropy} and Corollary~\ref{cor:one-zero} give
\begin{equation}\label{eq:one-sided-per}
 \per(B)\ge
 \max\{\kappa_n,\gamma_n\e^{s/2}\}.
\end{equation}
Consequently,
\begin{equation}\label{eq:one-sided-lower}
 \per(A)\ge
 (1-\tau)^n\max\{\kappa_n,\gamma_n\e^{s/2}\}.
\end{equation}
On the other hand, $\delta\ge\rho$, and therefore
\begin{equation}\label{eq:one-sided-upper}
 \per(A)=\gamma_n-\delta\le\gamma_n-\rho.
\end{equation}

For a one-sided pair, \eqref{eq:tau-bound}, $p\ge1$, and $\alpha_{n,b}\le n/2$ give
\begin{equation}\label{eq:one-sided-tau}
 \tau^2\le\frac{n\gamma_n}{2(1-\gamma_n)}=d_0^2<\frac1{100}.
\end{equation}

\subsection{The two entropy regimes}

Set
\[
 s_0=2\log\frac{\kappa_n}{\gamma_n}>0.
\]

\begin{lemma}[Small entropy]\label{lem:small-entropy}
If $s\le s_0$, then
\begin{equation}\label{eq:small-contradiction}
 \rho+\kappa_n(1-\tau)^n>\gamma_n.
\end{equation}
\end{lemma}

\begin{proof}
From \eqref{eq:rho-D-s},
\[
 \rho>Q_n\tau^2,
 \qquad
 Q_n=\frac{2(1-\gamma_n)^2}{5s_0}.
\]
Bernoulli's inequality gives
\[
 \rho+\kappa_n(1-\tau)^n-\gamma_n
 >\kappa_n-\gamma_n-n\kappa_n\tau+Q_n\tau^2.
\]
The right side is at least
\[
 \kappa_n-\gamma_n-\frac{n^2\kappa_n^2}{4Q_n}.
\]
Since $\log x\le x-1$,
\[
 s_0\le\frac{2(\kappa_n-\gamma_n)}{\gamma_n}.
\]
Thus the last expression is positive whenever
\begin{equation}\label{eq:small-certificate-condition}
 \frac25(1-\gamma_n)^2\gamma_n
 -\frac{n^2}{2}\kappa_n^2>0.
\end{equation}
The exact positive values of the left side are
\begin{equation}\label{eq:small-certificates}
\begin{array}{c|c}
 n&\text{value (decimal for scale)}\\ \hline
 8&7.682036\times10^{-4}\\
 9&3.378893\times10^{-4}\\
 10&1.383834\times10^{-4}.
\end{array}
\end{equation}
\end{proof}

\begin{lemma}[Large entropy]\label{lem:large-entropy}
If $s\ge s_0$, then
\begin{equation}\label{eq:large-contradiction}
 \rho+\gamma_n(1-\tau)^n\e^{s/2}>\gamma_n.
\end{equation}
\end{lemma}

\begin{proof}
Put $q=(1-\tau)^n$.  Since $\e^{s/2}\ge1+s/2$, equation \eqref{eq:rho-D-s} and AM--GM give
\[
 \begin{aligned}
 \rho+\gamma_nq\e^{s/2}-\gamma_n
 &>\frac25\frac{D^2}{s}
   +\gamma_nq\left(1+\frac{s}{2}\right)-\gamma_n\\
 &\ge D\sqrt{\frac45\gamma_nq}+\gamma_n(q-1).
 \end{aligned}
\]
By \eqref{eq:D-lower}, \eqref{eq:one-sided-tau}, and Bernoulli's inequality,
\[
 D\sqrt q
 \ge(1-\gamma_n)\tau\left(\frac9{10}\right)^{n/2},
 \qquad
 q-1\ge-n\tau.
\]
Therefore the preceding expression is at least
\[
 \tau\left[
 (1-\gamma_n)\left(\frac9{10}\right)^{n/2}
 \sqrt{\frac45\gamma_n}
 -n\gamma_n
 \right].
\]
The bracket is positive exactly when
\begin{equation}\label{eq:large-certificate-condition}
 \frac45(1-\gamma_n)^2\left(\frac9{10}\right)^n
 -n^2\gamma_n>0.
\end{equation}
The exact positive values are
\begin{equation}\label{eq:large-certificates}
\begin{array}{c|c}
 n&\text{value (decimal for scale)}\\ \hline
 8&1.889119\times10^{-1}\\
 9&2.334869\times10^{-1}\\
 10&2.424523\times10^{-1}.
\end{array}
\end{equation}
\end{proof}

Equations \eqref{eq:one-sided-lower} and \eqref{eq:one-sided-upper} contradict Lemma~\ref{lem:small-entropy} when $s\le s_0$ and Lemma~\ref{lem:large-entropy} when $s\ge s_0$.  Thus one-sided critical pairs are impossible.

\begin{theorem}[Boundary exclusion for $8\le n\le10$]\label{thm:boundary-low}
For every $n\in\{8,9,10\}$, no global maximizer of $\Phi$ on $K_n$ has a zero entry.
\end{theorem}

\begin{proof}
The case $\tau=0$ was excluded after Lemma~\ref{lem:hall-deficit}.  For $\tau>0$, Proposition~\ref{prop:low-two-sided} excludes two-sided critical pairs, and the one-sided analysis above excludes the remaining pairs.
\end{proof}

\section{Boundary exclusion in dimension seven}\label{sec:n7}

Dimension seven is the first dimension in which the scalar zero-block
comparison alone fails.  The failure is confined to critical Hall pairs with
$p=n-a-b=1$.  The complementary-minor entropy from
Theorem~\ref{thm:single-bridge-entropy} supplies exactly the missing
information.

Throughout this section let $n=7$, write
\[
 \gamma=\gamma_7=\frac{720}{117649},
\]
and let $A$ be a hypothetical boundary global maximizer.  The localization
parameter is
\[
 \eta_7=\frac1{47}.
\]
The two inequalities needed in Lemma~\ref{lem:subset-mid} are
\begin{equation}\label{eq:n7-localization}
 \eta_7^2-\frac{\gamma}{14(1-\gamma)}
 =\frac{23263}{1808073127}>0,
 \qquad
 \frac12-\frac37-\eta_7=\frac{33}{658}>0.
\end{equation}
Thus all subset-excess estimates of Section~\ref{sec:hall} apply in dimension
seven.  In particular,
\begin{equation}\label{eq:n7-alpha-max}
 \max_{1\le k\le6}\alpha_{7,k}
 =\alpha_{7,3}=\frac{53576}{15463}<\frac72.
\end{equation}
As before, the case $\tau=0$ is impossible, so fix a critical pair and use
$a,b,p,e,f,x$ as in Lemma~\ref{lem:hall-deficit}.

\subsection{Two-sided pairs with core size at least two}

Put
\begin{equation}\label{eq:n7-ordinary-parameters}
 T_7=\frac5{24},
 \qquad
 \lambda_7=\frac{329}{24}
 =\binom72-\binom73T_7.
\end{equation}
The common Hall estimate and \eqref{eq:n7-alpha-max} imply
\begin{equation}\label{eq:n7-ordinary-range}
 \tau^2\le\frac{7\gamma}{1-\gamma}<T_7^2,
 \qquad
 T_7^2-\frac{7\gamma}{1-\gamma}
 =\frac{20185}{67351104}>0.
\end{equation}
The alternating binomial argument of Lemma~\ref{lem:binomial-low} gives
\begin{equation}\label{eq:n7-ordinary-binomial}
 (1-t)^7\ge1-7t+\lambda_7t^2
 \qquad(0\le t\le T_7).
\end{equation}

\begin{proposition}[Ordinary dimension-seven margins]\label{prop:n7-ordinary}
For every $1\le a\le b$ with $p=7-a-b\ge2$,
\[
 M_7^{(2)}\left(L^{\sharp}_{7,a,b},
 \frac{\alpha_{7,a}+\alpha_{7,b}}{p^2}\right)>0,
\]
where \eqref{eq:M2} is evaluated with $\lambda=329/24$.
There are six pairs, and the minimum is attained at $(a,b)=(1,1)$:
\begin{equation}\label{eq:n7-ordinary-minimum}
 \frac{217499955444827479495}
 {21713873952987016007066496}>0
 \qquad(\approx1.001664\times10^{-5}).
\end{equation}
\end{proposition}

\begin{proof}
This is an exact rational substitution in Lemma~\ref{lem:second-order}.
The six comparisons are checked by the supplementary certificate described in
Appendix~\ref{app:code}.
\end{proof}

Proposition~\ref{prop:n7-ordinary} excludes every two-sided critical pair
with $p\ge2$.

\subsection{One-sided pairs}

The one-sided stationarity identity from
Lemma~\ref{lem:one-sided-stationarity} remains valid.  We sharpen only the
numerical energy estimates.  Define
\[
 d_7^2=\frac{7\gamma}{2(1-\gamma)}=\frac{2520}{116929}.
\]
Then
\begin{equation}\label{eq:n7-d7-gap}
 \frac9{400}-d_7^2=\frac{44361}{46771600}>0,
\end{equation}
so every row sum and every column sum of $A$ is greater than $17/20$.

\begin{lemma}[Dimension-seven product energy]\label{lem:n7-product-energy}
With $y_i=1-r_i^{-1}$, put
\[
 C_7=\frac{(1-\gamma)(17/20)^2}{2}.
\]
Then
\begin{equation}\label{eq:n7-product-energy}
 \rho\ge C_7\sum_i y_i^2,
 \qquad C_7>\frac13.
\end{equation}
Thus $\rho>\frac13\sum_i y_i^2$ unless all row sums equal one.  The
transposed statement holds for the column sums.
\end{lemma}

\begin{proof}
The proof of Lemma~\ref{lem:row-energy} gives the stated non-strict
bound with coefficient $C_7$.  This coefficient exceeds $1/3$, since
\[
 \frac{(1-\gamma)(17/20)^2}{2}-\frac13
 =\frac{7258243}{282357600}>0.
\]
\end{proof}

Retain the notation $D,s,q=(1-\tau)^7$ from the one-sided analysis in
Section~\ref{sec:low}.  Here $D>0$, so the reciprocal-deviation vector
$y$ is nonzero and the strict $1/3$ form of
Lemma~\ref{lem:n7-product-energy} applies.  Hence
\begin{equation}\label{eq:n7-one-sided-energy}
 \rho>\frac{D^2}{3s}
 \ge
 \frac{(1-\gamma)^2\tau^2}{3(1-\tau)^2s}.
\end{equation}
Also $\tau<3/20$ by \eqref{eq:n7-d7-gap}.

\begin{proposition}[One-sided exclusion in dimension seven]
\label{prop:n7-one-sided}
No one-sided critical pair occurs at a boundary global maximizer in dimension
seven.
\end{proposition}

\begin{proof}
Set $s_0=2\log(\kappa_7/\gamma)$.  If $s\le s_0$, the proof of
Lemma~\ref{lem:small-entropy}, with $1/3$ in place of $2/5$, reduces the
contradiction to
\begin{equation}\label{eq:n7-small-entropy-certificate}
 \frac13(1-\gamma)^2\gamma-\frac{49}{2}\kappa_7^2
 =
 \frac{22176489486262543850191055}
 {20672701805255574480144334848}>0.
\end{equation}

Suppose $s\ge s_0$.  From \eqref{eq:n7-one-sided-energy},
$\e^{s/2}\ge1+s/2$, and AM--GM,
\begin{equation}\label{eq:n7-large-start}
 \rho+\gamma q\e^{s/2}-\gamma
 >D\sqrt{\frac{2\gamma q}{3}}+\gamma(q-1).
\end{equation}
The stationarity bound gives
\[
 D\sqrt q\ge
 (1-\gamma)\tau(1-\tau)^{5/2}.
\]
Moreover, for $0\le t\le3/20$,
\[
 \frac{1-(1-t)^7}{t}
 \le7-21t+35t^2.
\]
The preceding estimates naturally reduce positivity to
\[
 \frac23(1-\gamma)^2(1-t)^5
 >\gamma\left(\frac{1-(1-t)^7}{t}\right)^2.
\]
Since $0\le t\le3/20$, one has $(1-t)^5\ge(1-t)^6$.  It is therefore
sufficient to verify the stronger inequality obtained by replacing the
left factor $(1-t)^5$ by the smaller $(1-t)^6$.  Together with the displayed
quadratic upper bound, this gives the sufficient condition
\begin{equation}\label{eq:n7-large-polynomial}
 W(t):=\frac23(1-\gamma)^2(1-t)^6
       -\gamma(7-21t+35t^2)^2>0.
\end{equation}
The degree-six Bernstein coefficients of $W$ on $[0,3/20]$ are all positive;
the smallest is
\begin{equation}\label{eq:n7-large-bernstein-min}
 \frac{155130361249919329}{1328763571296000000}>0.
\end{equation}
This exact coefficient calculation appears in the supplementary certificate.
Thus both entropy regimes contradict the upper bound
$\per(A)\le\gamma-\rho$, exactly as in Section~\ref{sec:low}.
\end{proof}

\subsection{Single-bridge pairs: preliminary reductions}

It remains to treat $p=1$.  By transposition assume $a\le b$, so
\begin{equation}\label{eq:n7-single-types}
 (a,b)\in\{(1,5),(2,4),(3,3)\}.
\end{equation}
Recall
\[
 e=\sum_{i\in I^c}r_i-a,
 \qquad
 f=\sum_{j\in J^c}c_j-b,
 \qquad
 x=s_A(I^c,J^c),
\]
so criticality gives
\begin{equation}\label{eq:n7-tau-efx}
 \tau=e+f-x.
\end{equation}

If $e\le0$, then $\tau\le f_+$ and the proof of
Lemma~\ref{lem:hall-deficit} gives the sharper estimate
\[
 \tau^2\le\alpha_{7,b}\frac{\delta}{1-\delta}.
\]
If $f\le0$, the analogous estimate uses $\alpha_{7,a}$.  Since
\begin{equation}\label{eq:n7-mixed-range}
 \left(\frac3{20}\right)^2
 -\alpha_{7,3}\frac{\gamma}{1-\gamma}
 =\frac{842770143}{723229250800}>0,
\end{equation}
we may use the quadratic binomial bound with
\[
 T=\frac3{20},
 \qquad
 \lambda=\binom72-\binom73T=\frac{63}{4}.
\]

\begin{proposition}[Sign-mixed single-bridge margins]
\label{prop:n7-sign-mixed}
For the three pairs in \eqref{eq:n7-single-types}, all six applications of
Lemma~\ref{lem:second-order} obtained from $e\le0$ or $f\le0$ have positive
margin.  Their minimum is
\begin{equation}\label{eq:n7-mixed-minimum}
 \frac{165035926385}{3897251304526692}>0
 \qquad(\approx4.234163\times10^{-5}),
\end{equation}
attained for $(a,b)=(1,5)$ with $e\le0$.
\end{proposition}

\begin{proof}
Substitute $L=L^{\sharp}_{7,a,b}$ and respectively
$c=\alpha_{7,b}$ or $c=\alpha_{7,a}$ into \eqref{eq:M2}, with
$\lambda=63/4$.  The supplementary script checks the six exact fractions.
\end{proof}

We may henceforth assume
\begin{equation}\label{eq:n7-positive-excesses}
 e>0,
 \qquad
 f>0.
\end{equation}
The subset estimates and \eqref{eq:n7-alpha-max} imply
\begin{equation}\label{eq:n7-ef-box}
 0<e,f<\frac3{20}.
\end{equation}
Put
\[
 t=e+f,
 \qquad
 m=7-a=b+1,
 \qquad
 \ell=7-b=a+1,
 \qquad
 d=\min\{a,b\}=a.
\]
Define the two group-product envelopes
\begin{equation}\label{eq:n7-product-envelopes}
 P_a(e)=\left(1+\frac ea\right)^a
        \left(1-\frac e{7-a}\right)^{7-a},
 \qquad
 P_b(f)=\left(1+\frac fb\right)^b
        \left(1-\frac f{7-b}\right)^{7-b}.
\end{equation}
If
\[
 P=\prod_i r_i,
 \qquad Q=\prod_jc_j,
 \qquad H=\per(A),
\]
write
\[
 u_r=P_a(e)-P,
 \qquad
 u_c=P_b(f)-Q,
\]
and set
\begin{equation}\label{eq:n7-DKu}
 D_0=2-P_a(e)-P_b(f),
 \qquad
 K=\gamma-D_0,
 \qquad
 u=\delta-D_0.
\end{equation}
Then
\begin{equation}\label{eq:n7-u-budget}
 u\ge u_r+u_c\ge0,
 \qquad
 H=K-u.
\end{equation}

\begin{lemma}[Conditional product energy]\label{lem:n7-conditional-energy}
Let
\[
 V_I=\sum_{i\in I}\left(
 \frac1{r_i}-\frac1m\sum_{h\in I}\frac1{r_h}
 \right)^2
\]
and define $V_J$ analogously from the columns indexed by $J$.  Then
\begin{equation}\label{eq:n7-conditional-energy}
 V_I\le3u_r,
 \qquad
 V_J\le3u_c.
\end{equation}
\end{lemma}

\begin{proof}
All row and column sums exceed $r_0=17/20$.  Let a group of $k$ positive
numbers $z_i$ have arithmetic mean $\mu$.  Since
$\sum_i(z_i/\mu-1)=0$,
\[
 k\log\mu-\sum_i\log z_i
 =\sum_i\left(\frac{z_i}{\mu}-1-
       \log\frac{z_i}{\mu}\right).
\]
For completeness, set $y=1-\mu/z_i$.  Then
\[
 \frac{z_i}{\mu}-1-\log\frac{z_i}{\mu}
 =\int_0^y\frac{t}{(1-t)^2}\,dt,
 \qquad
 \frac1{z_i}-\frac1\mu=-\frac{y}{\mu}.
\]
If $y\ge0$, the integral is at least $y^2/2$, while
$r_0^2/\mu^2\le1$.  If $y<0$, then on $[y,0]$,
\[
 (1-t)^{-2}\ge(1-y)^{-2}
 =\left(\frac{z_i}{\mu}\right)^2
 \ge\frac{r_0^2}{\mu^2}.
\]
Both cases give
\begin{equation}\label{eq:n7-group-log-energy}
 \frac{z_i}{\mu}-1-
 \log\frac{z_i}{\mu}
 \ge\frac{r_0^2}{2}
 \left(\frac1{z_i}-\frac1\mu\right)^2.
\end{equation}
Applying this to the two row groups and using the fact that variance is
minimized at the arithmetic mean of the reciprocals,
\[
 \log\frac{P_a(e)}P\ge\frac{r_0^2}{2}V_I.
\]
Therefore
\[
 u_r=P_a(e)-P
 =P\left(\exp\left(\log\frac{P_a(e)}P\right)-1\right)
 \ge(1-\gamma)\frac{r_0^2}{2}V_I\ge\frac13V_I,
\]
where the last inequality is the exact comparison in
Lemma~\ref{lem:n7-product-energy}.  The column proof is identical.
\end{proof}

\subsection{Core stationarity and the entropy squeeze}

Put $\theta=1-\tau$.  At a critical pair, the doubly stochastic matrix
$B\le A/\theta$ has $B[I^c,J^c]=0$.  The doubly stochastic cut identity
shows that $B[I,J]$ has total mass one.  Criticality and
\eqref{eq:n7-tau-efx} show that $(A/\theta)[I,J]$ has the same total mass.
Since the first matrix is entrywise dominated by the second, equality of
their total masses forces termwise saturation:
\begin{equation}\label{eq:n7-core-saturation}
 B[I,J]=\frac{A[I,J]}\theta.
\end{equation}
Let $\xi$ and $\zeta$ be the row- and column-sum vectors of this core and
put
\begin{equation}\label{eq:n7-core-s}
 s_I=\sum_{i\in I}\left(\xi_i-\frac1m\right)^2,
 \qquad
 s_J=\sum_{j\in J}\left(\zeta_j-\frac1\ell\right)^2,
 \qquad
 s=s_I+s_J.
\end{equation}
For $k\ge0$, let $H_k$ be the contribution to $H$ of permutations using
exactly $k$ entries of the block $I^c\times J^c$.  Put $H_0$ for the
zero-use contribution.

\begin{lemma}[Single-bridge stationarity]\label{lem:n7-core-stationarity}
One has
\begin{equation}\label{eq:n7-core-stationarity}
 P\left(\sum_{i\in I}\frac{\xi_i}{r_i}-1\right)
 +Q\left(\sum_{j\in J}\frac{\zeta_j}{c_j}-1\right)
 =\frac{\tau H+\sum_{k\ge1}kH_k}{1-\tau}.
\end{equation}
Moreover,
\begin{equation}\label{eq:n7-H0-entropy}
 H_0\ge
 L^{\sharp}_{7,a,b}(1-t)^7\e^{s/2}.
\end{equation}
\end{lemma}

\begin{proof}
Apply the common identity \eqref{eq:normalized-block-variation} to the core
$E=I\times J$.  Its total mass is $\theta$; by
\eqref{eq:n7-core-saturation}, its row and column masses are
$\theta\xi_i$ and $\theta\zeta_j$.  A permutation counted by $H_k$ uses
exactly $k+1$ core entries.  Substitution in
\eqref{eq:normalized-block-variation}, followed by division by $\theta$ and
use of $\Phi=P+Q-H$, gives exactly
\eqref{eq:n7-core-stationarity}.

For the second assertion, every permanent term of $B$ avoids the prescribed
zero block.  Entrywise domination $A\ge\theta B$ therefore maps each term of
$\per(B)$ to an $H_0$ term with at least $\theta^7$ times its weight.  Thus,
term by term,
\[
 H_0\ge \theta^7\per(B)
 \ge L^{\sharp}_{7,a,b}\theta^7\e^{s/2},
\]
where the last inequality is
Theorem~\ref{thm:single-bridge-entropy}.  Finally
$\theta=1-e-f+x\ge1-t$.
\end{proof}

For brevity write
\begin{equation}\label{eq:n7-A0}
 L=L^{\sharp}_{7,a,b},
 \qquad
 A_0=L(1-t)^7,
\end{equation}
and define
\begin{equation}\label{eq:n7-Bcal}
 \mathcal B=
 P_a(e)\frac{e}{m-e}
 +P_b(f)\frac{f}{\ell-f}
 -\frac{(t+d)K-dA_0}{1-t}.
\end{equation}

\begin{lemma}[Stationarity--entropy squeeze]\label{lem:n7-squeeze}
At a feasible single-bridge maximizer satisfying
\eqref{eq:n7-positive-excesses}, one has
\begin{equation}\label{eq:n7-entropy-lower}
 H\ge A_0\e^{s/2},
\end{equation}
and, whenever $\mathcal B>0$,
\begin{equation}\label{eq:n7-su-lower}
 su\ge\frac{\mathcal B^2}{3}.
\end{equation}
On the other hand,
\begin{equation}\label{eq:n7-su-upper}
 su\le\frac{(K-A_0)^2}{2A_0}.
\end{equation}
\end{lemma}

\begin{proof}
The first assertion is \eqref{eq:n7-H0-entropy} and $H\ge H_0$.
By the harmonic-mean inequality,
\[
 \sum_{i\in I}\frac{\xi_i}{r_i}-1
 \ge\frac{e}{m-e}-\sqrt{s_IV_I},
\]
and similarly in the column group.  Since a permutation uses at most $d$
entries of $I^c\times J^c$,
\[
 \sum_{k\ge1}kH_k\le d(H-H_0).
\]
The right side of \eqref{eq:n7-core-stationarity} is therefore at most
\begin{equation}\label{eq:n7-stationarity-upper}
 \frac{(\tau+d)(K-u)-dA_0}{1-\tau}
 \le
 \frac{(t+d)(K-u)-dA_0}{1-t}.
\end{equation}
Indeed, the derivative of the first rational function with respect to
$\tau$ is
\[
 \frac{(1+d)(K-u)-dA_0}{(1-\tau)^2}>0,
\]
because $K-u=H\ge H_0\ge A_0$, and then $\tau\le t$.

Put
\[
 E=\frac{e}{m-e},\qquad F=\frac{f}{\ell-f},\qquad
 c=\frac{t+d}{1-t},
\]
and denote the common value in \eqref{eq:n7-core-stationarity} by $Z$.
The harmonic-mean estimates and $P,Q\le1$ give
\begin{equation}\label{eq:n7-error-upper}
 PE+QF-Z\le\sqrt{s_IV_I}+\sqrt{s_JV_J}.
\end{equation}
On the other hand, \eqref{eq:n7-stationarity-upper} and
$P=P_a-u_r$, $Q=P_b-u_c$ give
\[
 PE+QF-Z
 \ge \mathcal B+cu-u_rE-u_cF.
\]
Here $c\ge\max\{E,F\}$, because
$E\le e/(1-e)\le t/(1-t)<c$, and similarly for $F$.
Together with $u\ge u_r+u_c$, this proves
$PE+QF-Z\ge\mathcal B$.  Lemma~\ref{lem:n7-conditional-energy} and
Cauchy--Schwarz now yield
\[
 \mathcal B
 \le\sqrt{s_IV_I}+\sqrt{s_JV_J}
 \le\sqrt{3s(u_r+u_c)}
 \le\sqrt{3su}.
\]
This proves \eqref{eq:n7-su-lower}.

Finally, $K-u=H\ge A_0\e^{s/2}$ and
$\e^{s/2}\ge1+s/2$ imply
\[
 su\le\frac{2u(K-A_0-u)}{A_0}
 \le\frac{(K-A_0)^2}{2A_0},
\]
which is \eqref{eq:n7-su-upper}.
\end{proof}

\subsection{The exact two-variable certificate}

\begin{lemma}[Exact Bernstein certificate principle]
\label{lem:bernstein-certificate}
Let $P$ be a rational polynomial on a rational box.  Express $P$ in the
tensor-product Bernstein basis of its coordinatewise degree.  If every
Bernstein coefficient is positive, then $P$ is positive on the box.
Moreover, exact de Casteljau subdivision at a rational parameter replaces a
parent box by finitely many child boxes whose union is the parent and whose
Bernstein coefficients are exact rational affine combinations of the parent
coefficients.  Consequently, a finite subdivision tree in which every leaf
has positive Bernstein coefficients is an exact certificate that $P>0$ on
the original box.
\end{lemma}

\begin{proof}
Each tensor-product Bernstein basis function is nonnegative on the box and
the basis functions sum to one, so the value of $P$ is a convex combination
of its Bernstein coefficients.  The de Casteljau identities are polynomial
identities and partition each coordinate interval exactly; iterating them
therefore covers the original box without gaps or overlaps except along
shared boundaries.  Rational subdivision points preserve exact rational
arithmetic.
\end{proof}

All quantities in the preceding subsection are rational functions of $e,f$.
The actual feasible set is contained in the larger closed square
\[
 \mathcal Q=\left[0,\frac3{20}\right]^2.
\]
We certify the entire square, so no separate semialgebraic description of
the transportation-feasible subset is required.  On $\mathcal Q$ one has
\[
 m-e>0,\qquad \ell-f>0,\qquad1-e-f\ge\frac7{10}>0,
\]
and also $A_0=L(1-e-f)^7>0$.  Hence
\[
 \Delta=(m-e)(\ell-f)(1-e-f)>0,
\]
so clearing denominators preserves every inequality direction.  Define the
three rational-coefficient polynomials
\[
 \mathcal G=A_0-K,
 \qquad
 \mathcal N=\Delta\mathcal B,
\]
and
\begin{equation}\label{eq:n7-GNF}
 \mathcal F=2A_0\mathcal N^2
 -3(K-A_0)^2\Delta^2.
\end{equation}
If $\mathcal G>0$, then \eqref{eq:n7-entropy-lower} contradicts
$H=K-u\le K$.  If $\mathcal N>0$ and $\mathcal F>0$, then
\eqref{eq:n7-su-lower} and \eqref{eq:n7-su-upper} contradict each other.
Thus it remains only to certify the following finite polynomial alternative.

\begin{proposition}[Exact single-bridge certificate]\label{prop:n7-bernstein}
For each $(a,b)\in\{(1,5),(2,4),(3,3)\}$ and every $(e,f)\in\mathcal Q$,
\begin{equation}\label{eq:n7-polynomial-alternative}
 \mathcal G(e,f)>0
 \quad\text{or}\quad
 \bigl(\mathcal N(e,f)>0\ \text{and}\ \mathcal F(e,f)>0\bigr).
\end{equation}
\end{proposition}

\begin{proof}
Apply Lemma~\ref{lem:bernstein-certificate}.  Starting from $\mathcal Q$,
the supplementary script subdivides the physically longer side at its
midpoint whenever neither branch in
\eqref{eq:n7-polynomial-alternative} is yet coefficientwise positive.  All
coefficients and all subdivision endpoints are instances of
\texttt{fractions.Fraction}; no floating-point operation or tolerance is
used.

The resulting finite trees are
\begin{equation}\label{eq:n7-bernstein-tree}
\begin{array}{c|rrrr}
(a,b)&\text{nodes}&\mathcal G\text{-leaves}&
 (\mathcal N,\mathcal F)\text{-leaves}&\text{maximum depth}\\ \hline
(1,5)&39&10&10&8\\
(2,4)&35&9&9&6\\
(3,3)&47&11&13&8.
\end{array}
\end{equation}
The complete coefficient generation and de Casteljau subdivision are given
in the supplementary certificate described in Appendix~\ref{app:code}.
\end{proof}

\begin{theorem}[Boundary exclusion in dimension seven]
\label{thm:boundary-seven}
No global maximizer of $\Phi$ on $K_7$ has a zero entry.
\end{theorem}

\begin{proof}
The case $\tau=0$ is excluded after Lemma~\ref{lem:hall-deficit}.
For $\tau>0$, Proposition~\ref{prop:n7-one-sided} excludes one-sided pairs,
Proposition~\ref{prop:n7-ordinary} excludes two-sided pairs with $p\ge2$,
and Proposition~\ref{prop:n7-sign-mixed} excludes single-bridge pairs with a
nonpositive complementary excess.  In the remaining case,
Lemma~\ref{lem:n7-squeeze} and Proposition~\ref{prop:n7-bernstein} give a
contradiction.  Hence no boundary global maximizer exists.
\end{proof}


\section{Boundary exclusion in dimension six}\label{sec:n6}

We now treat the remaining dimension.  Throughout this section,
\begin{equation}\label{eq:n6-gamma-kappa}
 n=6,\qquad
 \gamma=\gamma_6=\frac5{324},\qquad
 \kappa=\kappa_6=v_5G(5)=\frac{6144}{390625}>\gamma.
\end{equation}
Let $A$ be a hypothetical boundary global maximizer, and retain the notation
\[
 R=\prod_i r_i=1-\rho,\qquad
 C=\prod_jc_j=1-\sigma,\qquad
 H=\per(A)=\gamma-\delta
\]
from Section~\ref{sec:hall}.

\subsection{Uniform localization}

\begin{lemma}[Dimension-six localization]\label{lem:n6-localization}
For every nonempty proper set $E$ of rows, if
\[
 \sum_{i\in E}r_i=|E|+e,
\]
then
\begin{equation}\label{eq:n6-localization}
 |e|^2\le\frac{15}{319}
 <\left(\frac7{32}\right)^2
 <\left(\frac{11}{50}\right)^2.
\end{equation}
The same assertion holds for every column subset.  In particular,
\begin{equation}\label{eq:n6-singleton-lower}
 r_i,c_j>\frac{25}{32}\qquad(1\le i,j\le6).
\end{equation}
\end{lemma}

\begin{proof}
If $e>0$, separate AM--GM on $E$ and $E^c$, followed by binary Pinsker,
gives
\[
 \frac{2e^2}{6}\le-\log(1-\rho)
 \le\frac{\rho}{1-\rho}
 \le\frac{\gamma}{1-\gamma}.
\]
Thus $e^2\le15/319$.  If $e<0$, the complementary row set has positive
excess $-e$.  The column proof is identical.  The rational comparisons in
\eqref{eq:n6-localization} are verified exactly in the supplementary
certificate.
\end{proof}

Set
\begin{equation}\label{eq:n6-E0}
 E_0=\frac{11}{50}.
\end{equation}
As in the preceding dimensions, a boundary maximizer with Hall deficit
$\tau=0$ is already doubly stochastic and has a zero.  Corollary
\ref{cor:one-zero} would then give $H\ge\kappa>\gamma$, contradicting
Lemma~\ref{lem:joint-deficit}.  Hence $\tau>0$.

\subsection{Two-sided critical Hall pairs}

Assume first that a critical pair has $a,b\ge1$, and transpose so that
$a\le b$.  The six possible types are
\begin{equation}\label{eq:n6-types}
 (a,b,p)=(1,1,4),(1,2,3),(1,3,2),(2,2,2),(1,4,1),(2,3,1),
\end{equation}
where $p=6-a-b$.  Put
\begin{equation}\label{eq:n6-eftx}
 e=\sum_{i\in I^c}r_i-a,
 \qquad
 f=\sum_{j\in J^c}c_j-b,
 \qquad
 x=s_A(I^c,J^c),
 \qquad
 t=e+f.
\end{equation}
Criticality gives
\begin{equation}\label{eq:n6-tau-eft}
 p\tau=e+f-x=t-x.
\end{equation}
Consequently,
\begin{equation}\label{eq:n6-two-sided-domain}
 -E_0\le e,f\le E_0,
 \qquad
 0<t\le2E_0=\frac{11}{25}<p,
 \qquad
 \tau\le\frac tp.
\end{equation}
Put $\theta=1-\tau$.  The doubly stochastic cut identity and the forced
zero block give $s_B(I,J)=p$.  Criticality gives
$s_A(I,J)=p\theta$, so $B[I,J]\le(A/\theta)[I,J]$ and both sides have total
mass $p$.  Hence entrywise domination and equality of total masses imply the
termwise saturation
\[
 B[I,J]=\frac{A[I,J]}\theta.
\]
Let $m=6-a$, $\ell=6-b$, and let $\xi,\eta$ be the normalized row and
column marginals of this core, as in Proposition
\ref{prop:core-marginal-entropy}.  Put
\begin{equation}\label{eq:n6-Sxi}
 S_\xi=\sum_{i\in I}\left(\xi_i-\frac1m\right)^2,
 \qquad
 S_\eta=\sum_{j\in J}\left(\eta_j-\frac1\ell\right)^2.
\end{equation}

Decompose the permanent as
\[
 H=H_0+H_1+\cdots+H_d,
 \qquad d=\min(a,b)=a,
\]
where $H_k$ is the contribution of permutations using exactly $k$ entries
of $I^c\times J^c$.  Such a permutation uses exactly $p+k$ core entries.

\begin{lemma}[Dimension-six core stationarity]\label{lem:n6-core-stationarity}
At the global maximizer,
\begin{equation}\label{eq:n6-core-stationarity}
 R\left(\sum_{i\in I}\frac{\xi_i}{r_i}-1\right)
 +C\left(\sum_{j\in J}\frac{\eta_j}{c_j}-1\right)
 =\frac{\tau H+p^{-1}\sum_{k\ge1}kH_k}{1-\tau}.
\end{equation}
Consequently,
\begin{equation}\label{eq:n6-stationarity-upper}
 R\left(\sum_{i\in I}\frac{\xi_i}{r_i}-1\right)
 +C\left(\sum_{j\in J}\frac{\eta_j}{c_j}-1\right)
 \le\frac{(t+d)H-dH_0}{p-t}.
\end{equation}
\end{lemma}

\begin{proof}
Apply \eqref{eq:normalized-block-variation} to the core $I\times J$.  Its
mass is $p\theta$; its row and column masses are $p\theta\xi_i$ and
$p\theta\eta_j$; and a monomial in $H_k$ uses exactly $p+k$ core entries.
Substitution in the common identity, division by $p\theta$, and
$\Phi=R+C-H$ give \eqref{eq:n6-core-stationarity} term by term.

Since $k\le d$,
\[
 \sum_{k\ge1}kH_k\le d(H-H_0).
\]
Writing $z=p\tau$, the resulting upper bound is
\[
 f(z)=\frac{(z+d)H-dH_0}{p-z}.
\]
Its derivative has numerator
\[
 (p+d)H-dH_0\ge pH>0.
\]
Here $H>0$ because $A\ge\theta B$, $B$ is doubly stochastic, and hence
$H\ge\theta^6\per(B)\ge\theta^6\gamma_6>0$.  Thus $f$ is strictly
increasing.  Since $z=p\tau=t-x\le t$, substituting $t$ proves
\eqref{eq:n6-stationarity-upper}.
\end{proof}

Let $L=L^{\sharp}_{6,a,b}$.  Define $c_\xi,c_\eta$ by
\begin{equation}\label{eq:n6-entropy-coefficients}
 (c_\xi,c_\eta)=
 \begin{cases}
 \left(\dfrac{p^2}{2(a^2+b^2)},
       \dfrac{p^2}{2(a^2+b^2)}\right),&p\ge2,\\[3mm]
 \left(\dfrac12+\dfrac1{3m},
       \dfrac12+\dfrac1{3\ell}\right),&p=1.
 \end{cases}
\end{equation}
Put
\begin{equation}\label{eq:n6-A0}
 A_0=L\left(1-\frac tp\right)^6.
\end{equation}
Every permanent term of $B$ avoids the forced zero block.  Since
$A\ge\theta B$, termwise domination gives
$H_0\ge\theta^6\per(B)$.  Proposition
\ref{prop:core-marginal-entropy} and
$\theta=1-\tau\ge1-t/p$ therefore give
\begin{equation}\label{eq:n6-H0-lower}
 H_0\ge A_0(1+c_\xi S_\xi+c_\eta S_\eta).
\end{equation}
The joint deficit budget also gives
\begin{equation}\label{eq:n6-Hbar}
 H\le\overline H:=R+C-2+\gamma.
\end{equation}
It follows from Lemma~\ref{lem:n6-core-stationarity} that both
\begin{align}
 F={}&
 \frac{(t+d)\overline H-dA_0(1+c_\xi S_\xi+c_\eta S_\eta)}{p-t}
 \nonumber\\[-1mm]
 &-R\left(\sum_{i\in I}\frac{\xi_i}{r_i}-1\right)
 -C\left(\sum_{j\in J}\frac{\eta_j}{c_j}-1\right),\label{eq:n6-F}\\
 G={}&\overline H-A_0(1+c_\xi S_\xi+c_\eta S_\eta)\label{eq:n6-G}
\end{align}
are nonnegative.  Indeed, in \eqref{eq:n6-stationarity-upper} replace
$H$ by the larger quantity $\overline H$ and $H_0$ by the lower bound
\eqref{eq:n6-H0-lower}; this gives $F\ge0$.  Moreover
$\overline H\ge H\ge H_0$, which gives $G\ge0$.  Hence
$F+\zeta G\ge0$ for every $\zeta\ge0$.

\subsection{A two-level extremal reduction}

For a group of $r$ nonnegative numbers $z_i$ with total $s$, a probability
vector $w$, and constants $K,B_*,\lambda>0$, define on the closed product
of simplices the genuine polynomial
\begin{equation}\label{eq:n6-side-functional}
 \mathcal T(z,w)=
 KB_*\prod_{i=1}^r z_i
 -K\sum_{i=1}^r w_i\prod_{j\ne i}z_j
 -\lambda\sum_{i=1}^r\left(w_i-\frac1r\right)^2.
\end{equation}
When every $z_i>0$, this is equal to
\[
 K\prod_i z_i\left(B_*-\sum_i\frac{w_i}{z_i}\right)
 -\lambda\sum_i\left(w_i-\frac1r\right)^2.
\]

\begin{lemma}[Two-level side maximizer]\label{lem:n6-two-level}
Suppose $s/r>25/32$ and $B_*>32/25$.  The maximum of
\eqref{eq:n6-side-functional} is attained in the following form.  For some
$k\in\{1,\dots,r\}$,
\[
 w_i=\frac1k\ (1\le i\le k),
 \qquad
 w_i=0\ (i>k),
\]
and the $z_i$ have two levels,
\[
 z_i=h\ (i\le k),
 \qquad
 z_i=\ell_0\ (i>k),
 \qquad h\ge\ell_0.
\]
For $k<r$ they can be parametrized by $0\le u\le1$ as
\begin{equation}\label{eq:n6-two-level-param}
 h=\frac{s}{r}\left(1+u\frac{r-k}{k}\right),
 \qquad
 \ell_0=\frac{s}{r}(1-u).
\end{equation}
For $k=r$, one has $h=s/r$.
\end{lemma}

\begin{proof}
Compactness gives a maximizer.  If some $z_i=0$, the first term in the
polynomial form \eqref{eq:n6-side-functional} vanishes; the second and third
terms are nonpositive.  By contrast, at
$z_i=s/r$ and $w_i=1/r$ the value is
\[
 K(s/r)^r(B_*-r/s)>0,
\]
because $r/s<32/25<B_*$.  Hence every maximizing $z_i$ is positive, the
maximum is positive, and at a maximizer
$B_*-\sum_iw_i/z_i>0$.

Among all maximizers choose one minimizing
\[
 \mathcal V=\sum_i z_i^2+\sum_iw_i^2.
\]
Fix $i\ne j$ and preserve the two $z$- and $w$-sums by writing
\[
 z_i=h+x,\quad z_j=h-x,
 \qquad
 w_i=\omega+y,\quad w_j=\omega-y.
\]
Set
\[
 P_0=K\prod_{\nu\ne i,j}z_\nu>0,
 \qquad
 A=B_*-\sum_{\nu\ne i,j}\frac{w_\nu}{z_\nu}>0.
\]
Expanding the polynomial \eqref{eq:n6-side-functional}, all terms depending
on $(x,y)$ are, up to an additive constant,
\begin{equation}\label{eq:n6-pair-quadratic}
 -P_0A x^2+2P_0xy-2\lambda y^2.
\end{equation}
This explicit expansion is the reason for using the polynomial definition
on the closed domain.

Assume first that $w_i,w_j>0$.  Then $(x,y)$ lies in the relative interior
of the feasible pair rectangle, so both stationarity equations apply:
\[
 -2P_0Ax+2P_0y=0,
 \qquad
 2P_0x-4\lambda y=0.
\]
Either $x=y=0$, or $y=Ax$ and $P_0=2\lambda A$.  In the latter case
\eqref{eq:n6-pair-quadratic} is
\[
 -\frac{P_0}{A}(Ax-y)^2.
\]
For $0\le\theta\le1$, replacing $(x,y)$ by $(\theta x,\theta y)$ keeps
$z_i,z_j>0$ and $w_i,w_j\ge0$, preserves both pair sums, and remains on this
flat line.  It therefore preserves the maximum, whereas
\[
 \mathcal V(\theta)-\mathcal V(0)
 =2\theta^2(x^2+y^2)
\]
is strictly positive for $\theta>0$ unless $x=y=0$.  This contradicts the
minimal choice of $\mathcal V$.  Thus all positive coordinates of $w$ are
equal, and their corresponding $z_i$ are equal.

If $w_i=w_j=0$, then $y=0$ and the pair-dependent polynomial is the strictly
concave quadratic $-P_0Ax^2$.  Hence all inactive $z_i$ are equal.  Since
$w$ is a probability vector, if its support has size $k$, every active
coordinate equals $1/k$.

It remains to order the two $z$-levels.  Suppose $k<r$, transfer a small
amount of mass from one active $w$-coordinate to one inactive coordinate,
and keep $z$ fixed.  This is a feasible right-sided perturbation.  Its
directional derivative at zero is
\[
 K\prod_i z_i\left(\frac1h-\frac1{\ell_0}\right)
 +\frac{2\lambda}{k}\le0.
\]
If $h\le\ell_0$, the left side is strictly positive, a contradiction.
Thus $h>\ell_0$ when $k<r$; for $k=r$ there is only one level.  Solving the
fixed-total equation $kh+(r-k)\ell_0=s$ yields
\eqref{eq:n6-two-level-param}, with $0\le u\le1$.
\end{proof}

Set
\begin{equation}\label{eq:n6-zeta}
 \zeta=\begin{cases}14,&p\ge2,\\8,&p=1,\end{cases}
\end{equation}
and define
\begin{equation}\label{eq:n6-Bstar}
 B_*=\frac{p+d}{p-t}+\zeta,
 \qquad
 \lambda_\xi=c_\xi A_0\left(\frac d{p-t}+\zeta\right),
 \qquad
 \lambda_\eta=c_\eta A_0\left(\frac d{p-t}+\zeta\right).
\end{equation}
By \eqref{eq:n6-singleton-lower},
$\sum_i\xi_i/r_i<50/39<B_*$, and similarly for columns.  AM--GM on the
complementary groups therefore increases the upper bound for $F+\zeta G$
when $R,C$ are replaced by
\[
 \widehat R=\left(1+\frac ea\right)^a\prod_{i\in I}r_i,
 \qquad
 \widehat C=\left(1+\frac fb\right)^b\prod_{j\in J}c_j.
\]
Consequently,
\begin{equation}\label{eq:n6-dual-upper}
 F+\zeta G\le\mathcal U
 =C_0(t)+T_a(e,r,\xi)+T_b(f,c,\eta),
\end{equation}
where
\begin{align}
 C_0(t)&=
 \frac{(t+d)(\gamma-2)-dA_0}{p-t}
 +\zeta(\gamma-2-A_0),\label{eq:n6-C0}\\
 T_a&=
 \left(1+\frac ea\right)^a\prod_{i\in I}r_i
 \left(B_*-\sum_{i\in I}\frac{\xi_i}{r_i}\right)
 -\lambda_\xi S_\xi,\label{eq:n6-Ta}
\end{align}
and $T_b$ is the analogous column expression.

For a side of complement size $a$, group size $m=6-a$, support size $k$,
and excess $e$, Lemma~\ref{lem:n6-two-level} gives
\begin{equation}\label{eq:n6-Tcandidate}
 T_a\le
 \left(1+\frac ea\right)^a
 h^{k-1}\ell_0^{m-k}(B_*h-1)
 -\lambda_\xi\left(\frac1k-\frac1m\right),
\end{equation}
where $h,\ell_0$ are given by \eqref{eq:n6-two-level-param} with
$s=m-e$.  The formula for $k=m$ omits the low level and entropy term.  The
column side is identical.

\subsection{The exact four-variable certificate}

For fixed $0\le t\le2E_0$, the enlarged constraints
$-E_0\le e,f\le E_0$ and $e+f=t$ are equivalent to
$e\in[t-E_0,E_0]$ and $f=t-e$.  Thus the following map covers that entire
line segment exactly; adjoining the boundary $t=0$ only enlarges the actual
feasible set, which has $t>0$.  Map this closed triangular domain to the
unit square by
\begin{equation}\label{eq:n6-triangle-map}
 t=2E_0x,
 \qquad
 e=t-E_0+z(2E_0-t),
 \qquad
 f=E_0-z(2E_0-t),
 \qquad 0\le x,z\le1.
\end{equation}
For every type $(a,b)$, every pair of support sizes
\[
 1\le k\le6-a,
 \qquad
 1\le k'\le6-b,
\]
and the two spread parameters $u,v\in[0,1]$, substitute
\eqref{eq:n6-two-level-param}, \eqref{eq:n6-triangle-map}, and
\eqref{eq:n6-zeta} into the right side of
\eqref{eq:n6-dual-upper}.  Define
\begin{equation}\label{eq:n6-certificate-poly}
 P_{a,b,k,k'}(x,z,u,v)=-(p-t)\mathcal U.
\end{equation}
On the enlarged parameter cube,
\[
 p-t\ge p-\frac{11}{25}>0,
 \qquad A_0>0,
\]
and every row and column denominator is at least $25/32$.  Hence all
rational expressions used to define $\mathcal U$ have positive denominators,
and multiplication by $p-t$ reverses no inequality.  Thus this is a
rational polynomial.  Lemma~\ref{lem:bernstein-certificate}, implemented by
exact tensor-product Bernstein subdivision, proves
\begin{equation}\label{eq:n6-certificate-positive}
 P_{a,b,k,k'}>0\qquad\text{on }[0,1]^4
\end{equation}
for all $98$ cases.  The certificate summary is
\begin{equation}\label{eq:n6-two-sided-table}
\begin{array}{c@{\qquad}r@{\qquad}r@{\qquad}r}
\toprule
(a,b)&\text{polynomials}&\text{Bernstein nodes}&\text{maximum depth}\\
\midrule
(1,1)&25&165&6\\
(1,2)&20&276&7\\
(1,3)&15&483&14\\
(2,2)&16&416&13\\
(1,4)&10&352&14\\
(2,3)&12&558&15\\
\bottomrule
\end{array}
\end{equation}
Every terminal Bernstein coefficient is a strictly positive rational
number.  The smallest terminal coefficient for the tightest type $(2,3)$ is
\begin{equation}\label{eq:n6-smallest-fourvar}
 \frac{206812176444269171}
 {2388787200000000000000}>0
 \qquad(\approx8.657622\times10^{-5}).
\end{equation}
The supplementary script constructs every polynomial from
\eqref{eq:n6-certificate-poly}, performs exact power-to-Bernstein conversion,
and uses exact de Casteljau subdivision at $1/2$.

Since $p-t>0$, \eqref{eq:n6-certificate-positive} gives $\mathcal U<0$,
contradicting
\[
 0\le F+\zeta G\le\mathcal U.
\]
Thus no two-sided critical Hall pair can occur.

\subsection{One-sided critical Hall pairs}

After transposition, assume $I=[6]$.  As in Section~\ref{sec:low},
$b=0$ would force zero Hall deficit, so $1\le b\le5$ and
$|J|=p=6-b$.  The complementary $b$ columns have total mass $b+p\tau$;
applying Lemma~\ref{lem:n6-localization} to that positive excess gives
$(p\tau)^2\le15/319<E_0^2$.  Therefore
\begin{equation}\label{eq:n6-one-sided-mass}
 \sum_{j\in J}c_j=p(1-\tau),
 \qquad
 0<\tau<\frac{E_0}{p}.
\end{equation}
The corresponding columns of $B$ saturate $A/(1-\tau)$.  Put
\[
 m_i=\sum_{j\in J}b_{ij},
 \qquad
 \beta_i=\frac{m_i}{p},
 \qquad
 s=\sum_i\left(\beta_i-\frac16\right)^2.
\]
Lemma~\ref{lem:one-sided-stationarity} gives
\begin{equation}\label{eq:n6-D}
 D:=1-\sum_i\frac{\beta_i}{r_i}
 \ge\frac{(1-\gamma)\tau}{1-\tau}>0.
\end{equation}

\begin{lemma}[Dimension-six row energy]\label{lem:n6-row-energy}
With $y_i=1-r_i^{-1}$, put
\[
 C_6=\frac{(1-\gamma)(25/32)^2}{2}.
\]
Then
\begin{equation}\label{eq:n6-row-energy}
 \rho\ge C_6\sum_i y_i^2,
 \qquad C_6>\frac3{10}.
\end{equation}
In the present one-sided setting $D>0$, so $y\ne0$ and the comparison with
$3/10$ is strict.  Consequently,
\begin{equation}\label{eq:n6-rho-D}
 \rho>\frac3{10}
 \frac{(1-\gamma)^2\tau^2}{(1-\tau)^2s}.
\end{equation}
\end{lemma}

\begin{proof}
By Lemma~\ref{lem:n6-localization}, $r_i>25/32$.  The integral argument in
Lemma~\ref{lem:row-energy} gives
\[
 r-1-\log r\ge\frac{(25/32)^2}{2}
 \left(1-\frac1r\right)^2.
\]
Since $1-\rho\ge1-\gamma$,
\[
 \rho\ge(1-\rho)(-\log(1-\rho))
 \ge\frac{(1-\gamma)(25/32)^2}{2}\sum_i y_i^2.
\]
The exact gap
\begin{equation}\label{eq:n6-energy-gap}
 C_6-\frac3{10}
 =\frac{1547}{3317760}>0
\end{equation}
proves the stated coefficient comparison.  Since $D=\sum_i\beta_i y_i>0$,
$y$ is nonzero.  The harmonic-mean inequality gives $\sum_i y_i\le0$;
Cauchy--Schwarz together with \eqref{eq:n6-D} then proves
\eqref{eq:n6-rho-D}.  In particular, $s>0$.
\end{proof}

Retain the exact complementary column-product loss.  Two-group AM--GM gives
\begin{equation}\label{eq:n6-sigma-exact}
 \sigma\ge1-Q_b(\tau),
 \qquad
 Q_b(\tau)=
 \left(1+\frac{p\tau}{b}\right)^b(1-\tau)^p.
\end{equation}
Because $A$ is a boundary matrix, the dominated doubly stochastic matrix
$B$ has a zero.  If $q^{(j)}$ are the $p$ saturated columns, then convexity
of squared distance gives
\[
 s\le\frac1p\sum_{j\in J}s_6(q^{(j)})
 \le\max_{j\in J}s_6(q^{(j)}).
\]
Apply Theorem~\ref{thm:column-entropy} to the maximizing column and
Corollary~\ref{cor:one-zero}.  Since $A\ge(1-\tau)B$, we obtain
\begin{equation}\label{eq:n6-one-sided-per-lower}
 H\ge(1-\tau)^6\max\{\kappa,\gamma\exp(s/2)\}.
\end{equation}
On the other hand,
\begin{equation}\label{eq:n6-one-sided-per-upper}
 H\le\gamma-\rho-\sigma.
\end{equation}
Set $s_0=2\log(\kappa/\gamma)$.

If $s\le s_0$, then
$s_0\le2(\kappa-\gamma)/\gamma$, and
\eqref{eq:n6-rho-D} gives
\begin{equation}\label{eq:n6-Aenergy}
 \rho>A\frac{\tau^2}{(1-\tau)^2},
 \qquad
 A=\frac{3}{10}\frac{(1-\gamma)^2\gamma}{2(\kappa-\gamma)}
 =\frac{39750390625}{5253139008}.
\end{equation}
It is therefore enough to prove
\begin{equation}\label{eq:n6-small-poly}
 (1-\tau)^2
 \left[1-Q_b(\tau)+\kappa(1-\tau)^6-\gamma\right]
 +A\tau^2>0.
\end{equation}
For $b=1,\dots,5$, exact univariate Bernstein subdivision proves this on
$0\le\tau\le11/(50p)$; the required depths are respectively
$1,2,3,4,5$.

If $s\ge s_0$, use $\exp(s/2)\ge1+s/2$ and AM--GM in
\eqref{eq:n6-rho-D}.  With $q=1-\tau$,
\[
 \rho+\gamma q^6\exp(s/2)-\gamma
 >(1-\gamma)\tau q^2\sqrt{\frac{3\gamma}{5}}
   +\gamma(q^6-1).
\]
Here $3\gamma/5=1/108$, and
\begin{equation}\label{eq:n6-large-gap}
 \frac{(1-\gamma)^2}{108}-\frac9{1024}
 =\frac{33853}{181398528}>0.
\end{equation}
Thus it remains to prove, for $\tau>0$,
\begin{equation}\label{eq:n6-large-poly}
 1-Q_b(\tau)+\frac3{32}\tau(1-\tau)^2
 +\gamma\bigl((1-\tau)^6-1\bigr)>0.
\end{equation}
The left side has a factor $\tau$.  For all five values of $b$, the quotient
has globally positive Bernstein coefficients on $[0,11/(50p)]$; the
smallest coefficient is $1/864$.

Both entropy regimes contradict
\eqref{eq:n6-one-sided-per-lower}--\eqref{eq:n6-one-sided-per-upper}.
Hence no one-sided critical pair exists.

\begin{theorem}[Boundary exclusion in dimension six]
\label{thm:boundary-six}
No global maximizer of $\Phi$ on $K_6$ has a zero entry.
\end{theorem}

\begin{proof}
The case $\tau=0$ was excluded after Lemma~\ref{lem:n6-localization}.  If
$\tau>0$, a critical pair exists.  The preceding subsections exclude both
its two-sided and one-sided forms.  Thus no boundary global maximizer exists.
\end{proof}



\section{Proof of the main theorem and corollaries}\label{sec:completion}

\begin{proof}[Proof of Theorem~\ref{thm:main}]
Fix $6\le n\le15$.
\begin{enumerate}[label=\textup{(\arabic*)},leftmargin=2.4em]
\item Compactness of $K_n$ and continuity of $\Phi$ give a global maximizer
$A$.
\item Theorem~\ref{thm:boundary-six} for $n=6$,
Theorem~\ref{thm:boundary-seven} for $n=7$,
Theorem~\ref{thm:boundary-low} for $8\le n\le10$, and
Theorem~\ref{thm:boundary-mid-high} for $11\le n\le15$ exclude every
boundary maximizer.  Thus $A$ is positive.
\item The positive-maximizer theorem, Theorem~\ref{thm:positive-maximizer},
then gives $A=U_n$.
\item At $U_n$ all row and column sums equal one and
$\per(U_n)=n!/n^n$, so
\[
 \max_{A\in K_n}\Phi(A)=\Phi(U_n)=2-\frac{n!}{n^n}.
\]
Any matrix attaining this value is itself a global maximizer; repeating
steps (2)--(3) proves that it is $U_n$.  This establishes both the value and
uniqueness.
\end{enumerate}
\end{proof}

\begin{corollary}[No boundary maximizers]\label{cor:no-boundary}
For every $6\le n\le15$, every global maximizer of $\Phi$ lies in the relative interior of $K_n$.
\end{corollary}



\section{Discussion, scope, and verification status}\label{sec:discussion}

The proof has several logically separate modules.  Stable-polynomial
capacity first yields the support-sensitive permanent inequality, the sharp
rectangular-zero-block bound, and three entropy refinements.  The exact
complement entropy $\Psi_m$ controls both a selected supported column and the
weighted geometric mean of all complementary minors.  These two forms
combine into the core-marginal bound of Proposition
\ref{prop:core-marginal-entropy}.

The positive-maximizer theorem is independent of the boundary analysis.
Pairwise averaging is flat at a positive global maximizer, while the Hessian
at the uniform matrix is negative definite on the tangent hyperplane.
Cyclic averaging therefore identifies the unique interior maximizer.

The common boundary framework consists of the joint deficit budget,
relative-entropy subset estimates, Hall dilation, and support-sensitive
permanent lower bounds.  Scalar first- and second-order comparisons settle
dimensions eight through fifteen and all ordinary dimension-seven cuts.  The
three dimension-seven single-bridge types require complementary-minor
entropy, a core stationarity identity, and conditional product energy; the
last separation is the exact two-variable alternative in Proposition
\ref{prop:n7-bernstein}.

In dimension six, marginal entropy must be retained even when the Hall core
has size $p\ge2$.  The exact stationarity identity in Lemma
\ref{lem:n6-core-stationarity} and the two-level reduction in Lemma
\ref{lem:n6-two-level} turn every two-sided configuration into one of $98$
four-variable polynomial inequalities.  The one-sided argument is sharper
than its higher-dimensional counterpart because it keeps the full
complementary column-product envelope $Q_b(\tau)$.

The computer-assisted component consists of finite exact rational
certificates: scalar comparisons in all dimensions, the dimension-seven
one- and two-variable Bernstein certificates, the $98$ dimension-six
four-variable certificates, and ten dimension-six univariate certificates.
The analytic reductions, the polynomials being certified, the positivity
criterion, and all subdivision statistics are stated in the text.  No
floating-point comparison, tolerance, random choice, or numerical optimizer
is used.

The theorem established here is exactly the range $6\le n\le15$.  We do not
infer a complete classification of the remaining dimensions from the
historical references: the cited $n=4$ paper establishes full
indecomposability of maximizers, while the $n=16$ and $n\ge17$ claims are
recent preprints as of the version date.  None of these external dimension
claims is used in the proof above.

\paragraph{Meaning of ``unconditional'' and verification status.}
The proof does not assume any previously claimed dimension of Dittert's
conjecture, nor any other unproved conjecture.  It does use standard results,
among them compactness and continuity, weighted AM--GM, Cauchy--Schwarz,
Bernoulli's inequality, convexity, the harmonic-mean inequality,
max-flow/min-cut, Hurwitz's theorem, Gauss--Lucas, and finite-dimensional
Taylor expansion.  The finite comparisons are exact rational computations;
the author checkers and independent reconstructions are included, and no
success condition is implemented by a Python \texttt{assert}.  A successful
run therefore performs the same checks under ordinary and optimized Python
execution.  This is an exact computer-assisted proof in the usual
mathematical sense, but it is not a Lean, Coq, Isabelle, or other
proof-assistant-kernel formalization.  Human refereeing of the analytic
reductions remains essential.

\section*{Declarations}
\paragraph{Code availability.}
All generators, exact checkers, independent audit implementations, tests,
source hashes, and the machine-readable verification manifest are contained
in the supplementary archive.  The accepted version should additionally
cite an immutable public repository release and archival DOI; those external
identifiers cannot be assigned inside an anonymous review package.

\paragraph{Data availability.}
No external data set is used.  Every finite datum required by the proof is
generated from the displayed rational formulas and is included in the
supplementary package.

\paragraph{Funding and conflicts of interest.}
No funding or conflict information is supplied in the anonymous version.
The corresponding author must replace this sentence with the applicable
journal declarations before acceptance.

\paragraph{Author contributions.}
The author(s) are responsible for the mathematical argument, the exact
certificate software, and the final verification.  A journal-specific CRediT
statement should be inserted when the author list is disclosed.

\paragraph{Use of computational and language-assistance tools.}
Computer algebra and exact Python programs were used for formula expansion,
certificate generation, checking, and manuscript preparation.  Language and
software-assistance tools were also used during revision.  All mathematical
claims, source transformations, and final outputs were reviewed and accepted
by the author(s), who retain full responsibility for the work.

\appendix

\section{Exact rational certificates}\label{app:certificates}

This appendix records representative exact comparisons used in the text.
The supplementary programs verify every scalar margin and every Bernstein
subdivision.

For dimension six,
\[
 \gamma_6=\frac5{324},\qquad
 \kappa_6=\frac{6144}{390625},\qquad
 \frac{15}{319}<\left(\frac7{32}\right)^2
 <\left(\frac{11}{50}\right)^2.
\]
The row-energy and large-entropy gaps are
\[
 \frac{(1-\gamma_6)(25/32)^2}{2}-\frac3{10}
 =\frac{1547}{3317760},
 \qquad
 \frac{(1-\gamma_6)^2}{108}-\frac9{1024}
 =\frac{33853}{181398528}.
\]
The four-variable certificate is summarized in
\eqref{eq:n6-two-sided-table}, and the smallest recorded terminal
coefficient for its tightest type is
\eqref{eq:n6-smallest-fourvar}.  The one-sided certificates consist of five
small-entropy polynomials and five large-entropy quotients; the latter have
minimum Bernstein coefficient $1/864$.

For dimension seven, the localization and half-gap comparisons are
\begin{equation}\label{eq:app-n7-localization}
 \eta_7^2-\frac{\gamma_7}{14(1-\gamma_7)}
 =\frac{23263}{1808073127},
 \qquad
 \frac12-\frac37-\eta_7=\frac{33}{658}.
\end{equation}
The minimum ordinary two-sided margin is
\[
 \frac{217499955444827479495}
 {21713873952987016007066496},
\]
and the minimum sign-mixed single-bridge margin is
\[
 \frac{165035926385}{3897251304526692}.
\]
The one-sided energy, small-entropy, and large-entropy certificates are the
positive quantities in \eqref{eq:n7-product-energy},
\eqref{eq:n7-small-entropy-certificate}, and
\eqref{eq:n7-large-bernstein-min}.  The exact two-variable certificate is
summarized in \eqref{eq:n7-bernstein-tree}.

The remaining exact margins are displayed in equations
\eqref{eq:mid-margin-table}, \eqref{eq:n11-cstar-gap},
\eqref{eq:n11-special-margin}, \eqref{eq:low-two-sided-table},
\eqref{eq:small-certificates}, and \eqref{eq:large-certificates}.

For the localization constants in \eqref{eq:eta-table}, one has
\begin{equation}\label{eq:localization-certificates}
\begin{array}{c|c|c}
 n&\displaystyle
 \eta_n^2-\frac{\gamma_n}{2n(1-\gamma_n)}
 &\displaystyle
 \frac12-\frac{\lfloor(n-1)/2\rfloor}{n}-\eta_n\\ \hline
 7&\dfrac{23263}{1808073127}&\dfrac{33}{658}\\[3mm]
 8&\dfrac{4757}{836844800}&\dfrac9{80}\\[3mm]
 9&\dfrac{242489}{87087962025}&\dfrac{13}{270}\\[3mm]
 10&\dfrac{31109}{41313127850}&\dfrac{11}{115}\\[3mm]
 11&\dfrac{63007753811}{34945789841847500}&\dfrac{82}{1925}\\[3mm]
 12&\dfrac{11651173}{90828753405000}&\dfrac{319}{3900}\\[3mm]
 13&\dfrac{13072911571453}{366471344281458130000}&\dfrac{537}{14300}\\[3mm]
 14&\dfrac{71438752181}{2511343447066440000}&\dfrac{893}{12600}.
\end{array}
\end{equation}
Every entry is positive.

The auxiliary comparison used for $11\le n\le14$ is
\begin{equation}\label{eq:gamma11-small}
 \frac1{1000}-\gamma_{11}
 =\frac{22308624601}{25937424601000}>0.
\end{equation}
For $n=8,9,10$, the inequalities $\gamma_n<1/81$ follow from
\begin{equation}\label{eq:gamma-low-small}
\begin{array}{c|c}
 n&1/81-\gamma_n\\ \hline
 8&\dfrac{105557}{10616832}\\[2mm]
 9&\dfrac{54569}{4782969}\\[2mm]
 10&\dfrac{1516573}{126562500}.
\end{array}
\end{equation}

For dimension $15$, the six inequalities in \eqref{eq:n15-gamma-bounds}--\eqref{eq:n15-mu2-bounds} are equivalent to the positivity of the following integers:
\begin{equation}\label{eq:n15-integers}
\begin{aligned}
15^{15}-300000\cdot15!
 &=45\,591\,579\,980\,859\,375,\\
29863\cdot15^{15}-10^{10}15!
 &=81\,568\,443\,603\,515\,625,\\
10^{10}13!13^{13}-29937\,14^{26}
 &=290\,584\,966\,207\,055\,614\,607\,247\,802\,368,\\
3\,14^{26}-10^6 13!13^{13}
 &=3\,939\,835\,293\,455\,120\,947\,239\,518\,208,\\
10^8 13!12^{24}-301\,13^{37}
 &=18\,198\,910\,159\,813\,162\,803\,331\,251\,792\,701\,065\,161\,767,\\
4\,13^{37}-10^6 13!12^{24}
 &=162\,574\,843\,672\,041\,436\,633\,515\,011\,954\,046\,933\,353\,332.
\end{aligned}
\end{equation}

\section{Exact-arithmetic verification}\label{app:code}

The supplementary package contains two author verifiers, three
independent reconstruction checkers, exact unit tests, and one orchestration
script.  The principal executable files are listed below.
\begin{center}
\small
\begin{tabularx}{\textwidth}{@{}>{\ttfamily}l Y@{}}
\toprule
\normalfont File & Purpose \\
\midrule
dittert\_n6\_certificates.py
 & dimension-six generator and exact certificates; \\
dittert\_n7\_n15\_certificates.py
 & standard-library certificates for dimensions seven through fifteen; \\
independent\_n6\_fast\_audit.py
 & independent reconstruction of all $98$ four-variable certificates; \\
independent\_n6\_univariate\_audit.py
 & independent standard-library reconstruction of the ten one-variable certificates; \\
independent\_n7\_certificate\_audit.py
 & independent symbolic and Bernstein reconstruction of the three single-bridge trees; \\
run\_all\_certificates.py
 & environment-aware orchestration, hashing, manifest generation, and regression checks. \\
\bottomrule
\end{tabularx}
\end{center}
Run
\[
 \texttt{python3 run\_all\_certificates.py}.
\]
A successful full run ends with
\begin{center}
\texttt{All Dittert certificates for dimensions 6 through 15 passed.}
\end{center}
It also writes the machine-readable files
\begin{center}
\path{verification_manifest.json}\quad and\quad
\path{verification_output.txt}.
\end{center}
Success is printed only after every subprocess exits normally and all
expected statistics have been checked.

The dimension-seven-through-fifteen verifier requires Python~3.11 or newer
and only the standard library.  It evaluates every localization, Hall-range,
first-order, and second-order margin as a \texttt{Fraction}; constructs the
univariate and bivariate dimension-seven polynomials; performs exact
power-to-Bernstein conversion and exact dyadic de Casteljau subdivision; and
checks all reported tree statistics literally.

The dimension-six verifier was tested with Python~3.11--3.13; the archived
reference environment uses Python~3.13, SymPy~1.14, and NumPy~2.3.  SymPy is used only to expand the explicit rational polynomials,
and NumPy only stores object arrays.  Every generated coefficient is
converted to \texttt{fractions.Fraction} before Bernstein conversion,
subdivision, or comparison.  The program constructs all $98$ polynomials in
\eqref{eq:n6-certificate-poly}, checks the node and depth counts in
\eqref{eq:n6-two-sided-table}, and verifies the ten one-sided polynomials.
The accompanying \texttt{requirements.txt} records the tested dependencies.

\begin{table}[ht]
\centering
\small
\begin{tabularx}{\textwidth}{@{}Y Y Y r@{}}
\toprule
Textual object & Author-code entry point & Enumeration & Count \\
\midrule
Dimension-six two-sided alternative, \eqref{eq:n6-certificate-poly}
 & \texttt{two\_sided\_family}, \texttt{verify\_two\_sided}
 & six Hall types and all row/column support pairs & $98$ \\
Dimension-six one-sided bounds
 & \texttt{verify\_one\_sided}
 & $b=1,\ldots,5$, two entropy regimes & $10$ \\
Dimension-seven single-bridge alternative, Proposition~\ref{prop:n7-bernstein}
 & \path{single_bridge_polynomials}, \path{certify_single_bridge}
 & $(a,b)=(1,5),(2,4),(3,3)$ & $3$ trees \\
Dimensions $8$--$10$ second-order margins
 & \texttt{verify\_n8\_n10}
 & all admissible two-sided Hall types & $12+16+20$ \\
Dimensions $11$--$14$ first-order margins
 & \texttt{verify\_n11\_n14}
 & all listed Hall types, with one $n=11$ refinement & $35+42+49+56$ \\
Dimension $15$ thin--thick bounds
 & \texttt{verify\_n15}
 & two rational margins and six cleared-denominator tests & $8$ \\
\bottomrule
\end{tabularx}
\caption{Map from the analytic reductions to the exact verification code.}
\label{tab:formula-code-map}
\end{table}

Neither verifier contains a floating-point comparison, tolerance, random
choice, numerical optimization routine, or validation implemented by a
Python \texttt{assert}.  Ordinary and optimized Python execution therefore
perform the same mathematical checks.  The Bernstein implication is
elementary: on a box, tensor-product Bernstein basis functions are
nonnegative and sum to one, so a polynomial lies between the smallest and
largest of its Bernstein coefficients.

\begin{thebibliography}{99}

\bibitem{CheonWanless2012}
G.-S. Cheon and I. M. Wanless,
\emph{Some results towards the Dittert conjecture on permanents},
Linear Algebra Appl. \textbf{436} (2012), 791--801.

\bibitem{Egorychev1981}
G. P. Egorychev,
\emph{The solution of van der Waerden's problem for permanents},
Soviet Math. Dokl. \textbf{23} (1981), 619--622.

\bibitem{Falikman1981}
D. I. Falikman,
\emph{Proof of the van der Waerden conjecture regarding the permanent of a doubly stochastic matrix},
Math. Notes Acad. Sci. USSR \textbf{29} (1981), 475--479.

\bibitem{Gurvits2008}
L. Gurvits,
\emph{Van der Waerden/Schrijver--Valiant like conjectures and stable (aka hyperbolic) homogeneous polynomials: one theorem for all},
Electron. J. Combin. \textbf{15} (2008), Research Paper 66.

\bibitem{Hwang1986}
S.-G. Hwang,
\emph{A note on a conjecture on permanents},
Linear Algebra Appl. \textbf{76} (1986), 31--44.

\bibitem{Hwang1987}
S.-G. Hwang,
\emph{On a conjecture of E. Dittert},
Linear Algebra Appl. \textbf{95} (1987), 161--169.

\bibitem{Kafidov2026}
B. Kafidov,
\emph{Dittert's conjecture in dimension 16 via a joint-deficit scaling lemma},
arXiv:2607.19439, preprint, version cited 25 July 2026.

\bibitem{KnoppSinkhorn1982}
P. Knopp and R. Sinkhorn,
\emph{Minimum permanents of doubly stochastic matrices with at least one zero entry},
Linear Multilinear Algebra \textbf{11} (1982), 351--355.

\bibitem{Minc1983}
H. Minc,
\emph{Theory of permanents 1978--1981},
Linear Multilinear Algebra \textbf{12} (1983), 227--263.

\bibitem{Pang2026}
Z. Pang,
\emph{Proof of Dittert's conjecture for dimensions $n\ge17$},
arXiv:2606.01531, preprint, version cited 25 July 2026.

\bibitem{Sinkhorn1984}
R. Sinkhorn,
\emph{A problem related to the van der Waerden permanent theorem},
Linear Multilinear Algebra \textbf{16} (1984), 167--173.

\bibitem{UdayanSomasundaram2024}
D. K. Udayan and K. Somasundaram,
\emph{Lih Wang's and Dittert's conjectures on permanents},
Special Matrices \textbf{12} (2024), Article 20240006,\
\url{https://doi.org/10.1515/spma-2024-0006}.

\bibitem{Zhan2007}
X. Zhan,
\emph{Open problems in matrix theory},
in: Proceedings of the International Congress of Chinese Mathematicians 2007, Vol.~II, 1--17.

\end{thebibliography}

\end{document}
